\documentclass[12pt,a4paper]{article}
\usepackage[onehalfspacing]{setspace}
\usepackage{amssymb,amsmath,amsthm,amsfonts,geometry,graphicx,caption,natbib,booktabs,xcolor,bm,float,array}
\usepackage{lmodern}
\usepackage{xr-hyper}
\usepackage[colorlinks=true,linkcolor=cyan,urlcolor=cyan,citecolor=cyan,hypertexnames=false]{hyperref}
\externaldocument{supplement}
% Exhibits are read from the results folders of the two packages and from figures/.
\makeatletter
\def\input@path{{figures/}{../results/paper/}{../simulations/results/mc_diff/}{../simulations/results/tex/}{../simulations/results/figures/}}
\makeatother
\graphicspath{{figures/}{../results/paper/}{../simulations/results/figures/}}
\geometry{left=1.25in,right=1.25in,top=1.25in,bottom=1.25in}
\setlength{\emergencystretch}{3em}

\newtheorem{thm}{Theorem}
\newtheorem{prop}{Proposition}
\newtheorem{lem}{Lemma}
\theoremstyle{definition}
\newtheorem{assumption}{Assumption}
\newtheorem{remark}{Remark}
\newtheorem{assumptionb}{Assumption}\renewcommand{\theassumptionb}{3$'$}
\newtheorem{assumptionc}{Assumption}\renewcommand{\theassumptionc}{4$'$}
\theoremstyle{definition}\newtheorem{example}{Example}
\theoremstyle{definition}\newtheorem{definition}{Definition}

\newcommand{\todo}[1]{{\color{red}[\textbf{TODO:} #1]}}
\newcommand{\safegraphic}[2][]{\IfFileExists{#2}{\includegraphics[#1]{#2}}{\fbox{\todo{figure file \detokenize{#2} missing}}}}
\newcommand{\safeinput}[1]{\IfFileExists{#1}{\input{#1}}{\fbox{\todo{table file \detokenize{#1} missing}}}}
\newcommand{\E}{\mathbb{E}}
\newcommand{\cov}{\operatorname{Cov}}
\newcommand{\var}{\operatorname{Var}}

\title{Control VAR: a counterfactual based approach to inference in macroeconomics}
\author{Raimondo Pala\thanks{University of Rome Tor Vergata. Email: raimondopala@gmail.com}}
\date{\today}

\begin{document}
\maketitle

\begin{abstract}
\noindent A recursive VAR is read as causal whenever the policy is ordered first. I show that this ordering is not sufficient when the policy reacts to the economy it is meant to affect: the impulse response then combines the causal effect with a bias from the common trends that the policy reacts to and a bias specific to the treated unit. I develop the Control VAR (CVAR), which compares the treated unit with untreated controls that share its stochastic trends, and removes the first bias and, under exogeneity with respect to the counterfactual cycle, the second. The method allows repeated, continuous and multiple policies, separates the effect of a single policy change from that of its later values, measures the bias of a VAR without controls through the response of the controls, and bounds the remaining bias. I estimate the weights, derive bootstrap inference, show that the CVAR in changes remains valid when the treated unit has a stochastic trend of its own, test whether the more efficient CVAR in levels is supported, and characterise the bias of weights estimated in changes. I apply the CVAR to the Alaska Permanent Fund Dividend, extended unemployment benefits and California's Low Carbon Fuel Standard.

\bigskip
\noindent\textbf{JEL Classification:} C32, C22, C21, H24, J22, J65, Q48.\\
\noindent\textbf{Keywords:} Vector autoregressions, causal inference, control units, cointegration, cash transfers, unemployment insurance, fuel standards.
\end{abstract}

\newpage
\section{Introduction\label{sec:intro}}

What would have happened to inflation had central banks around the world not introduced unconventional monetary policy at the onset of the Great Financial Crisis? Is monetary policy effective in controlling inflation? What is the impact of monetary policy on the economy?

While some of these questions may appear too tough to tackle, economists have them frequently in mind, and often find themselves at a loss as to which methods to use to answer them. Long-run, short-run and sign restrictions in VARs, local projections and linear regressions are just some of the many tools developed to address them, yet economists still face the difficult reality of being unable to give a properly causal answer. Many of these questions naturally involve the formulation of a counterfactual. When we think about the effectiveness of unconventional monetary policy, we have in mind a world in which central banks preferred inaction in the face of a looming disinflationary pressure, and when we think about the effects of a natural disaster, we have in mind a world in which no disaster happened. We rarely make the case, however, that we are truly identifying the effect of the policy over inaction, or the effect of the disaster, and for a long time causality has been almost a forbidden topic in macroeconomics.

The difficulty is specific to the social sciences, where the lack of experiments leaves economists to make informed guesses. \citet{Nakamura2018a} compare economists to weather forecasters, painting a gloomy picture of economic science, but they also express optimism about the direction the field has taken. Much of the recent literature has developed tools that formulate a counterfactual in the potential outcomes framework of \citet{Rubin1980} for VARs and local projections \citep{bojinovshephard2019,RambachanSheppard2021,dube2023local}, and another stream defines counterfactuals through optimal policies \citep{mckay2023can}. Both define the effect of a macroeconomic policy as the difference between two alternative assignments of the policy.

The tool that macroeconomists most often use to measure such effects is the recursive vector autoregression, in which the policy is ordered before the outcomes and its unexpected movements are treated as shocks whose impulse responses give the effect of the policy \citep{Sims1980,ChristianoEichenbaumEvans1999,Bloom2009}. The causal literature has given this reading a foundation, since the impulse response is a causal effect when the policy does not react, within the period, to the shocks that move the outcome \citep{RambachanSheppard2021}. Many policies fail this requirement by design, because they are meant to react to the economy: central banks set rates in response to current inflation and activity, fiscal stabilisers expand when output falls, and unemployment benefits are extended when a state's unemployment rises. A countercyclical policy is therefore applied when the economy is at its weakest, so the recovery that would have followed anyway is attributed to the policy, which appears to work even when it has no effect. I show that ordering such a policy first is not sufficient for a causal reading of its impulse response, and I give the exact form of the bias.

Applied microeconomics faces treatments that depend on the outcome in many settings, and it has well-established answers when units that do not receive the treatment are available, such as difference-in-differences \citep{Card1990,CardKrueger,Abadie2005} and synthetic control \citep{AbadieGardeazabal2003,Abadie2010}. Two features of macroeconomic data limit their use. Comparable units are few, and macroeconomic series wander without returning to a fixed mean, so that regressions among them can produce spurious relations \citep{GrangerNewbold1974}. Cointegration showed that such series can share a common stochastic trend, so that a combination of them is stable \citep{EngleGranger1987,Johansen1991}, and synthetic control rests on the same idea when it reconstructs the path a treated unit would have followed from a weighted combination of untreated units that share its trends \citep{Abadie2015,HsiaoChingWan2012,Arkhangelsky2021}.

This paper brings the two traditions together in the control VAR (CVAR), which uses control units to remove the common trends of the treated unit and estimates a VAR in the policy and the remaining cycle. Since the common trends carry the shocks the policy reacts to, removing them removes the confounding they create. Theorem \ref{thm:var} makes the confounding precise for a recursive VAR without controls, whose impulse response combines a causal component with two biases. The causal component is a weighted average of effects whose form depends on the distribution of the policy, an average treatment effect on the treated for a dummy and an average causal response for a continuous policy. The first bias comes from the common trends the policy reacts to, so that an intervention that only partly offset a recession suffered also by economies that did not intervene can appear harmful, and the second comes from the shocks that only the treated unit experiences. Replacing the outcome with its cycle removes the first bias (Theorem \ref{thm:cycle}), and when the policy is exogenous with respect to the cycle the second vanishes, so that the response of the cycle is the effect of a policy surprise (Theorem \ref{thm:ident}). That response includes the later reaction of the policy to its own consequences, which is the object macroeconomists usually report, and Theorem \ref{thm:direct} separates it from the effect of a single change in the policy with its later values held fixed, the impulse causal effect of \citet{RambachanSheppard2021}.

The controls also give users of VARs a direct measure of the bias they remove. When the weighted controls are added to a VAR, their response to a policy that never reached them is the bias from the common trends (Proposition \ref{prop:placebo}), and the same response bounds the bias specific to the treated unit when exogeneity is in doubt, giving the multiple of it at which a conclusion breaks down (Proposition \ref{prop:bounds}). Because the weights on the controls are estimated, Theorem \ref{thm:est} and Proposition \ref{prop:J} show how their estimation enters the distribution of the response and when it can be ignored.

The controls remove the common trends only if they share every stochastic trend of the treated unit. Over several decades a region can also follow a path of its own, through the rise or decline of an industry or a population that grows faster than that of its neighbours, and comparisons in levels, synthetic control among them, assume such a path away. The CVAR does not need this assumption, because estimated on the changes of the cycle it identifies the same effects when the treated unit carries a trend of its own (Theorem \ref{thm:diff}), and I use it as the default. The CVAR in levels is more efficient when the treated unit shares all its trends with the controls, and Theorem \ref{thm:po} gives a test of the null hypothesis that it does not, so that a rejection allows the estimator in levels while a non-rejection only costs efficiency. Applied to the employment, income, house prices and consumer prices of US states and regions, each compared with its neighbours, the test rejects about as often as its nominal size, which shows that the assumption behind comparisons in levels cannot be supported for most units even over several decades.

The results rest on five assumptions: the consistency condition of the potential outcomes framework for time series \citep{RambachanSheppard2021} together with a stable unit treatment value assumption \citep{Rubin1980}, under which the policy does not reach the controls; the factor structure of synthetic control \citep{Abadie2010}; an effect linear and additive in the policy, as in the literature on VARs \citep{KilianLutkepohl2017}; a conditional exogeneity condition of the kind used for policy evaluation in time series \citep{AngristJordaKuersteiner2018}, imposed on the cycle instead of the outcome; and stationarity. The Supplementary Appendix relaxes the first and the third, recovering the effect when the policy spills over to the controls through a known exposure mapping and when the effect is not linear in the surprise. Its simulations, calibrated to the application to unemployment benefits, show that a VAR without controls overstates the effect of a policy that reacts to common shocks while the CVAR is unbiased, and that with long enough lags a bootstrap that re-estimates the weights in every draw delivers bands with close to nominal coverage.

I apply the CVAR to three policies that react to the economy they affect. The dividend that Alaska pays every year to its residents follows the returns of a sovereign wealth fund, so its largest payments came after the stock market booms that preceded the recessions of 2001 and 2008. A VAR on Alaska alone consequently finds that a larger dividend lowers employment, a response that the states that never receive the dividend also show, while the CVAR finds no fall in employment and a rise in leisure and hospitality, the sector where the spending of the dividend lands \citep{Kueng2018}. Where synthetic control has evaluated the introduction of the dividend \citep{JonesMarinescu2022}, the CVAR uses the variation in its size over more than thirty payments. Extended unemployment benefits switch on in recessions, so a VAR without controls finds that they sharply lower unemployment, and the states that never extended benefits show almost the same response; the CVAR finds effects close to zero, in line with the evidence for periods in which benefit durations were already long \citep{AcostaEtAl2025}. Finally, the cost of the Low Carbon Fuel Standard of California changes every month with the price of its credits, and the CVAR finds a pass-through to the retail price of gasoline consistent with full pass-through, with narrower bands than a VAR without controls.

The CVAR belongs to the methods that build counterfactuals for a treated time series from untreated units. The artificial counterfactual \citep{CarvalhoMasiniMedeiros2018,MasiniMedeiros2021} and CausalImpact \citep{Brodersen2015} target the effect of a single intervention, and synthetic control \citep{Abadie2010} and synthetic difference-in-differences \citep{Arkhangelsky2021} estimate the weights from the periods before a single adoption. Difference-in-differences with staggered adoption \citep{callawaysantannamultipletime2021,sunabraham2021} and with non-binary and non-absorbing treatments \citep{callawaysantannabacon2021,deChaisemartinDHaultfoeuille2020} identifies the effect when the timing of the treatment is unrelated to the outcomes once unit and time effects are removed. Local-projection difference-in-differences \citep{dube2023local} and panel VARs with a causal interpretation \citep{pala2024pvarcontrol} allow a recurring policy that reacts to the past, although they are restricted to dummy policies and do not account for cointegration among the variables. The CVAR covers the remaining case of a policy that recurs, takes several values and reacts to shocks it shares with the controls, and because it estimates the weights from the stochastic trends over the whole sample it needs no period free of the policy. It thereby relaxes the central condition under which the causal interpretation of VARs and local projections holds \citep{RambachanSheppard2021,GoncalvesHerreraKilianPesavento2021,kolesar2024dynamic}. Table \ref{tab:methods} summarises what each approach allows.

\begin{table}[!t]
\centering
\caption{What each approach allows\label{tab:methods}}
\footnotesize
\setlength{\tabcolsep}{4pt}
\begin{tabular}{>{\raggedright\arraybackslash}p{4.6cm}ccccc}
\toprule
 & Recurring & Continuous & Reaction to & Heterogeneous & Stochastic \\
 & policy & intensity & the past & exposure to & trends \\
 & & & & common shocks & \\
\midrule
DiD, staggered adoption & $\times$ & $\times$ & $\times$ & $\times$ & $\times$ \\
DiD, non-binary and non-absorbing & \checkmark & \checkmark & $\times$ & $\times$ & $\times$ \\
Local-projection DiD & \checkmark & $\times$ & \checkmark & $\times$ & $\times$ \\
Synthetic control, synthetic DiD & $\times$ & $\times$ & $\times$ & \checkmark & $\times$ \\
Artificial counterfactual, CausalImpact & $\times$ & $\times$ & $\times$ & \checkmark & \checkmark \\
Recursive VARs, local projections & \checkmark & \checkmark & \checkmark & $\times$ & -- \\
Panel VARs & \checkmark & $\times$ & \checkmark & $\times$ & $\times$ \\
\midrule
CVAR & \checkmark & \checkmark & \checkmark & \checkmark & \checkmark \\
\bottomrule
\end{tabular}
\par\smallskip\parbox{0.95\textwidth}{\footnotesize DiD: difference-in-differences. Recurring policy: the policy switches on and off or changes repeatedly. Continuous intensity: the policy takes more than two values. Reaction to the past: the policy may respond to the past of the outcome it affects. Heterogeneous exposure to common shocks: the treated and untreated units may respond differently to the shocks they share, and the policy may react to those shocks. Stochastic trends: the outcomes may share stochastic trends, with cointegration among the units. Recursive VARs and local projections use no untreated units, so the last column does not apply to them. Difference-in-differences, staggered adoption: \citet{callawaysantannamultipletime2021,sunabraham2021}; non-binary and non-absorbing: \citet{callawaysantannabacon2021,deChaisemartinDHaultfoeuille2020}; local projections: \citet{dube2023local}; synthetic control: \citet{Abadie2010,Arkhangelsky2021}; artificial counterfactual: \citet{CarvalhoMasiniMedeiros2018,MasiniMedeiros2021,Brodersen2015}; VARs and local projections: \citet{RambachanSheppard2021,kolesar2024dynamic}; panel VARs: \citet{pala2024pvarcontrol}.}
\end{table}

Section \ref{sec:causal} presents the causal framework and Section \ref{sec:assumptions} the assumptions. Section \ref{sec:theorems} contains the theorems, estimation and inference, Section \ref{sec:test} the test of the design, Section \ref{sec:ext} the case of a treated unit with a trend of its own, Sections \ref{sec:alaska} to \ref{sec:lcfssum} the applications, and Section \ref{sec:conclusions} concludes.

\section{Causal framework\label{sec:causal}}

\citet{RambachanSheppard2021} define the dynamic causal effects of a policy in a time series through potential outcomes indexed by the whole path of the policy. I follow their example and adapt the framework to a setting with one treated unit and a set of control units, building on the panel formulation of \citet{bojinovshephard2019} and on the potential outcomes of \citet{Rubin1980}.

There are $K+1$ units: a treated unit, indexed by $1$, and $K$ control units, indexed by $k=1,\dots,K$. For $t=1,\dots,T$, each unit $i$ receives a policy $W_{it}$, and a vector $R_t$ of other policies may also be in place. In the application, the units are states, $W_{it}$ is the dividend that state $i$ pays to each of its residents in month $t$, in thousands of dollars, which is positive only in Alaska and only in the month of the yearly payment, and $R_t$ is the monthly change in the price of oil. The policy may also be a dummy, such as an indicator for whether a programme is in place.

Let $\mathbf{w}_{1:T}=(w_{i,1:T})_{i}$ collect a path of the policy for every unit over the whole sample, and let $r_{1:T}$ be a path of the other policies. The potential outcome of unit $i$ in period $t$ is
\begin{equation}
Y_{it}(\mathbf{w}_{1:T},r_{1:T}),
\label{eq:po}
\end{equation}
the value its outcome would take if the policies of all units had followed those paths. In the application, $Y_{1t}(\mathbf{w}_{1:T},r_{1:T})$ is employment in Alaska in month $t$ under a given history of dividends. In this general form, employment in Alaska today may depend on the whole path of the policies: on past dividends, on the expectation of future ones, and on the other policies, such as the current price of oil. The observed outcomes are the potential outcomes at the realised paths, $Y_{it}=Y_{it}(\mathbf{W}_{1:T},R_{1:T})$.

Causal effects are comparisons of potential outcomes across paths. The most general is the effect of a whole path, $Y_{1,t+h}(\mathbf{w}_{1:T},r_{1:T})-Y_{1,t+h}(\mathbf{w}'_{1:T},r_{1:T})$. Its realised counterpart compares the realised path with a path in which the treated unit never receives the policy,
\begin{equation}
\Delta_{t+h}=Y_{1,t+h}(\mathbf{W}_{1:T},R_{1:T})-Y_{1,t+h}(\mathbf{W}^{0}_{1:T},R_{1:T}),
\label{eq:effectpath}
\end{equation}
where $\mathbf{W}^0_{1:T}$ replaces the path of the treated unit with zeros and leaves the paths of the other units unchanged. It is the effect on employment in Alaska of all the dividends that Alaska actually paid up to $t+h$.

A policy that reacts to the economy is itself an outcome of its history, and a change in the policy today changes the policy tomorrow. I therefore also treat future policies as potential outcomes. For a period $t$, let $W_{1,t+k}(w)$, $k\geq0$, be the policy of the treated unit in period $t+k$ if its policy in period $t$ is set to $w$, the history before $t$ is the realised one, and the policy then follows its usual determinants, so that $W_{1t}(w)=w$ and the realised policy is $W_{1,t+k}=W_{1,t+k}(W_{1t})$. This is the formulation of dynamic treatment regimes \citep{Robins1986,Murphy2003}, in which the treatment in each period responds to the history that earlier treatments helped to create.

Two causal effects of the policy in period $t$ follow from this formulation, and they differ in what happens to the policy after period $t$. The first, introduced by \citet{RambachanSheppard2021}, changes the policy in period $t$ and holds the rest of the path fixed.

\begin{definition}[Impulse causal effect\label{def:ice}]
The impulse causal effect at horizon $h$ of setting the policy of the treated unit in period $t$ to $w$ instead of $w'$ is
\begin{equation}
Y_{1,t+h}(\mathbf{W}_{1:t-1},w,\mathbf{W}_{t+1:T},R_{1:T})-Y_{1,t+h}(\mathbf{W}_{1:t-1},w',\mathbf{W}_{t+1:T},R_{1:T}),
\label{eq:ice}
\end{equation}
with the policies of every unit in every other period at their realised values.
\end{definition}

It measures the direct effect of changing the policy in period $t$ on the outcome $h$ periods later. The second lets the policy follow its usual determinants after period $t$.

\begin{definition}[Regime effect\label{def:regime}]
The regime effect at horizon $h$ of setting the policy of the treated unit in period $t$ to $w$ instead of $w'$ is
\begin{equation}
\begin{split}
&Y_{1,t+h}\langle w\rangle-Y_{1,t+h}\langle w'\rangle,\quad\text{where}\\
&Y_{1,t+h}\langle w\rangle=Y_{1,t+h}\big(\mathbf{W}_{1:t-1},w,W_{1,t+1}(w),\dots,W_{1,t+h}(w),\mathbf{W}_{-1,t:T},R_{1:T}\big),
\end{split}
\label{eq:regime}
\end{equation}
where $\mathbf{W}_{-1,t:T}$ are the realised paths of the other units.
\end{definition}

It measures the effect of changing the policy in period $t$ together with the changes in its later values that the usual determinants of the policy then bring about. A policymaker who considers a different policy today, knowing that the policy will then evolve as usual, is interested in the regime effect.

A recursive VAR reports a third object, which connects $\Delta_{t+h}$ to the timing of the policy through its unexpected part. Let $Z_t=(R_t',W_{1t},G_t)'$, where $G_t$ is the cycle of the treated unit, let $\mathcal{F}_{t-1}$ be the information in $Z_{t-1},Z_{t-2},\dots$, and let
\begin{equation}
u_t=W_{1t}-\E^\ast[W_{1t}\mid R_t,\mathcal{F}_{t-1}]
\label{eq:innovation}
\end{equation}
be the innovation of the policy, where $\E^\ast$ denotes the linear projection. In the application, $u_t$ is the part of this month's dividend that could not be predicted from the past dividends, the price of oil and the past of the economy of Alaska. The third causal effect weighs the effect of the realised path by this innovation.

\begin{definition}[Effect of a policy surprise\label{def:surprise}]
The effect of a policy surprise at horizon $h$ is
\begin{equation}
\theta_h=\frac{\E[\Delta_{t+h}\,u_t]}{\E[u_t^2]}.
\label{eq:estimand}
\end{equation}
\end{definition}

In essence, $\theta_h$ is the part of the effect of the dividends on employment $h$ months later that is associated with an unexpected dividend of one unit in month $t$. Like the regime effect, it includes the effect of the dividends that follow the surprise. 

\section{Assumptions\label{sec:assumptions}}

The first assumption restricts how the policies enter the potential outcomes. Its first two parts are needed by any time-series estimate of the effect of a policy, with or without controls, and its third part is specific to the use of controls.

\begin{assumption}[No anticipation, no interference and untreated controls\label{assu:A1}]
\begin{enumerate}
\item[(i)] No anticipation: $Y_{it}(\mathbf{w}_{1:T},r_{1:T})=Y_{it}(\mathbf{w}_{1:t},r_{1:t})$ for every unit $i$ and period $t$.
\item[(ii)] No interference: $Y_{1t}(\mathbf{w}_{1:T},r_{1:T})=Y_{1t}(w_{1,1:T},r_{1:T})$, and $Y_{kt}(\mathbf{w}_{1:T},r_{1:T})=Y_{kt}(\mathbf{w}'_{1:T},r_{1:T})$ for every control $k$ and every pair of paths $\mathbf{w}_{1:T}$, $\mathbf{w}'_{1:T}$ that differ only in the path of the treated unit.
\item[(iii)] Untreated controls: $W_{kt}=0$ for every control $k$ and period $t$.
\end{enumerate}
\end{assumption}

Part (i) is the no anticipation condition of \citet{RambachanSheppard2021}: a policy affects the outcome only once it is in place, so that, in the example of the dividend, payments that have not yet been made do not affect employment in Alaska. Part (ii) is the stable unit treatment value assumption of \citet{Rubin1980}, adapted to time series. As in any comparison of treated and untreated units, it makes the units comparable by ruling out spillovers: the policy of the treated unit does not change the potential outcomes of the controls. In the example, it rules out an effect of the dividend on employment in other states, through the spending of Alaskans elsewhere or through movements of people into or out of Alaska.\footnote{Supplementary Appendix \ref{sec:spillnl} relaxes this part when the policy spills over to the controls through a known exposure.} Part (iii) requires that the controls never receive the policy, which holds by construction when they are states that pay no dividend.

Under Assumption \ref{assu:A1} the notation simplifies. I write $W_t$ for the policy of the treated unit, $Y^1_t(w_{-\infty:t})$ for the potential outcome of the treated unit under its own path, holding the other policies at their realised values, $Y^1_t(0)$ for its value had it never received the policy, and $Y^0_{kt}$ for the observed outcomes of the controls, which are their potential outcomes without the policy. The effect of the policy actually applied in \eqref{eq:effectpath} becomes $\Delta_{t+h}=Y^1_{t+h}-Y^1_{t+h}(0)$.

\begin{assumption}[Common trends\label{assu:A2}]
There is an $r$-dimensional process $F_t$, integrated of order one with stationary first difference, and a $K\times r$ matrix $\Lambda$ of rank $r\leq K$, such that
\begin{equation}
Y^0_t=\Lambda F_t+v_t,\qquad Y^1_t(0)=\lambda'F_t+e_t,
\label{eq:factors}
\end{equation}
where $v_t$ and $e_t$ are stationary and $\lambda=\Lambda'b$ for some $b\in\mathbb{R}^K$.
\end{assumption}

Assumption \ref{assu:A2} is a factor model: the outcome of the treated unit, such as employment in Alaska, follows stochastic trends that it shares with the outcomes of the controls, such as employment in other states. The vector $F_t$ collects the forces behind these trends, such as the national business cycle or the price of oil, and each unit loads on them differently. The condition $\lambda=\Lambda'b$ says that the treated unit loads on these forces as a weighted combination of the controls does. The model resembles the factor model of synthetic control \citep{Abadie2010}, from which it differs in the nature of the trends, which are integrated as in the cointegration literature \citep{EngleGranger1987}, and in the role of the weights. Under it, the counterfactual cycle of the treated unit,
\begin{equation}
G^0_t=Y^1_t(0)-b'Y^0_t=e_t-b'v_t,
\label{eq:cycle0}
\end{equation}
is stationary, with mean $a$. The weights $b$ play the role of the synthetic control weights, and I call $C_t=b'Y^0_t$ the weighted controls, so that the observed cycle is $G_t=Y^1_t-C_t$.

Two features of $b$ differ from the usual practice of synthetic control, and both follow from the role of the weights here. The weights are not required to sum to one. Employment in Alaska may be more or less cyclical than in the controls, and the weights must reproduce the amplitude of its exposure to the trends, in addition to its direction. The weights are also not required to be unique. When there are more controls than common trends, $r<K$, the controls are cointegrated among themselves, the combinations $S_t=H'Y^0_t=H'v_t$ are stationary for any $K\times(K-r)$ matrix $H$ whose columns span the null space of $\Lambda'$, and every vector in $\mathcal{B}=\{b+H\phi:\phi\in\mathbb{R}^{K-r}\}$ removes the trends. This is the situation of synthetic control with more donors than factors \citep{AbadieLHour2021}. The economist should therefore pick the controls that can best approximate movements in the treated economy.

\begin{assumption}[Additive dynamic effects\label{assu:A3}]
$Y^1_t(w_{-\infty:t})=Y^1_t(0)+\sum_{j\geq0}\tau_jw_{t-j}$, with $\sum_{j\geq0}|\tau_j|<\infty$.
\end{assumption}

Assumption \ref{assu:A3} is the linearity that the literature on VARs imposes when it interprets impulse responses \citep{KilianLutkepohl2017}. Supplementary Appendix \ref{sec:identification} gives the weights with which the response averages heterogeneous effects when the effect is not linear, and Supplementary Appendix \ref{sec:spillnl} shows how to recover the response to a surprise of a given size and sign.

Under Assumption \ref{assu:A3}, the impulse causal effect \eqref{eq:ice} of a unit change in the policy in period $t$ on the outcome $h$ periods later is $\tau_h$, and the effect of a policy surprise in \eqref{eq:estimand} decomposes as
\begin{equation}
\theta_h=\sum_{j=0}^h\tau_j\,\psi_{h-j},\qquad \psi_k=\frac{\E[W_{t+k}u_t]}{\E[u_t^2]},
\label{eq:decomp}
\end{equation}
where $\psi_k$ is the path that the policy is expected to follow $k$ periods after a surprise.\footnote{The decomposition uses the orthogonality of $u_t$ to the past of the policy, so that only the policy from period $t$ onward contributes.} The regime effect requires an assumption on the potential policies.

\begin{assumptionb}[Linear policy response\label{assu:A3b}]
$W_{t+k}(w)=W_{t+k}(0)+\phi_kw$ for $k\geq0$, with $\phi_0=1$ and $\sum_{k\geq0}|\phi_k|<\infty$.
\end{assumptionb}

Assumption \ref{assu:A3b} says that a change of $w$ in the policy today changes its value $k$ periods later by $\phi_kw$, whatever path it would otherwise have followed, so that $\phi_k$ is the propagation of the policy through its own determinants. Under Assumptions \ref{assu:A3} and \ref{assu:A3b}, the regime effect \eqref{eq:regime} of a unit change in the policy in period $t$ is $\theta^\ast_h=\sum_{j=0}^h\tau_j\phi_{h-j}$. It combines the direct effects of the policy with its own propagation, and it is the object that macroeconomists usually report as the response of the outcome to a policy shock. The effect of a surprise satisfies
\begin{equation}
\theta_h=\theta^\ast_h+\sum_{j=0}^{h-1}\tau_j\,\kappa_{h-j},\qquad \kappa_k=\frac{\E[W_{t+k}(0)\,u_t]}{\E[u_t^2]}.
\label{eq:pathbias}
\end{equation}
Since $\psi_k=\phi_k+\kappa_k$, the coefficient $\kappa_k$ is the part of the path of the policy $k$ periods ahead that a surprise predicts without causing it. It is non-zero when the surprise carries news about later values of the policy that would have occurred anyway, and the effect of a surprise then attributes to the policy the effect of those later values. When $\kappa_k=0$ for every $k$, the surprise moves the later policy only through its propagation, and the effect of a surprise equals the regime effect.

Under Assumptions \ref{assu:A1} to \ref{assu:A3}, the observed cycle $G_t=Y^1_t-b'Y^0_t$ satisfies
\begin{equation}
G_t=G^0_t+\sum_{j\geq0}\tau_jW_{t-j}.
\label{eq:cycle}
\end{equation}

\begin{assumption}[Exogeneity with respect to the counterfactual cycle\label{assu:A4}]
$\E[G^0_{t+h}u_t]=0$ for all $h\geq0$.
\end{assumption}

Assumption \ref{assu:A4} is a conditional exogeneity condition of the kind used to evaluate policies in time series \citep{AngristJordaKuersteiner2018}, imposed on the counterfactual cycle in place of the outcome. It allows the policy to react to the past of the treated unit and of the controls, to the other policies, and to the contemporaneous innovations of the common trends. It requires that, given the past, the policy surprise does not predict the part of the outcome of the treated unit that the controls do not share and that would have occurred without the policy; in the example, a larger dividend must not predict a movement of employment specific to Alaska.

\begin{assumption}[Stationarity and moments\label{assu:A5}]
$(R_t',W_t,G^0_t,\Delta F_t',v_t')'$ is covariance stationary with finite fourth moments, and $\E[u_t^2]>0$.
\end{assumption}

Assumption \ref{assu:A5} requires that the policies, the counterfactual cycle and the growth of the common trends have stable means and variances.

\section{Theorems\label{sec:theorems}}

\subsection{What a VAR without controls estimates\label{sec:var}}

Before turning to the CVAR, I show what a standard recursive VAR without controls estimates. Let $\tilde u_t$ be the innovation of the policy in a recursive VAR in $(R_t',W_t,Y^1_t)'$, with the past of $Y^1_t$ among the regressors. The impulse response of the outcome of the treated unit $h$ periods after a surprise is then $\gamma^{V}_h=\E[(Y^1_{t+h}-Y^1_{t-1})\tilde u_t]/\E[\tilde u_t^2]$, and $\tilde\theta_h$ denotes the effect of a policy surprise \eqref{eq:estimand} computed with $\tilde u_t$.\footnote{By Lemma \ref{lem:outcome} in Supplementary Appendix \ref{sec:identification}, $\gamma^V_h$ is the response of a VAR in levels, or of a VAR in differences cumulated to levels, when the lags of the system span the predictable part of the policy.}

\begin{thm}[What a VAR without controls identifies\label{thm:var}]
Under Assumptions \ref{assu:A1}, \ref{assu:A2}, \ref{assu:A3} and \ref{assu:A5},
\begin{equation}
\gamma^{V}_h=\tilde\theta_h+\lambda'\frac{\E[(F_{t+h}-F_{t-1})\tilde u_t]}{\E[\tilde u_t^2]}+\frac{\E[(e_{t+h}-e_{t-1})\tilde u_t]}{\E[\tilde u_t^2]}.
\label{eq:bias}
\end{equation}
\end{thm}

The response of the VAR has three parts. The first, $\tilde\theta_h$, is the causal component: the effect on the treated unit of a surprise in the policy, which is the object the researcher wants to estimate.\footnote{When the effect of the policy is not linear, the causal component is a weighted average of effects, with weights that depend on the distribution of the policy. Supplementary Appendix \ref{sec:identification} characterises it for dummy, continuous and non-negative policies, with results close to those of \citet{kolesar2024dynamic} for local projections and VARs.} The second is the bias from the common trends: the average movement of the common trends from the period before a surprise to $h$ periods after it, per unit of surprise, which $\lambda$ translates into the outcome of the treated unit. In the example of the dividend, the common trends include the national business cycle and the financial markets whose returns set the size of the dividend, and the VAR attributes to the dividend the movements of these trends that follow it. The common trends are not the only source of bias. The third term depends on $e_t$, the part of the outcome of the treated unit that neither the common trends nor the policy explain. It captures the reverse causality that arises when the policy reacts to developments specific to the treated unit, and it is zero when the policy surprise is unrelated to the later movements of $e_t$, the exogeneity under which \citet{RambachanSheppard2021} give the VAR its causal interpretation.

\subsection{The CVAR\label{sec:cvar}}

The CVAR subtracts the weighted controls from the outcome of the treated unit. Because the controls share the common trends, the subtraction removes the first source of bias, the common developments that the policy reacts to, and leaves the bias specific to the treated unit, now measured on the counterfactual cycle, as the next theorem shows.

\begin{thm}[What the CVAR identifies\label{thm:cycle}]
Under Assumptions \ref{assu:A1}, \ref{assu:A2}, \ref{assu:A3} and \ref{assu:A5}, the response of the cycle to the policy innovation is
\begin{equation}
\gamma_h\equiv\frac{\E[G_{t+h}u_t]}{\E[u_t^2]}=\theta_h+\frac{\E[G^0_{t+h}u_t]}{\E[u_t^2]},\qquad h=0,1,2,\dots
\label{eq:cyclebias}
\end{equation}
If $Z_t$ follows a VAR, possibly of infinite order, $\gamma_h$ is the response of $G_t$ at horizon $h$ to a unit innovation in $W_t$ in the recursive identification with $R_t$ first, $W_t$ second and $G_t$ last.
\end{thm}

Comparing Theorem \ref{thm:cycle} with Theorem \ref{thm:var}, the common-trend term has disappeared, because the weighted controls load on the common trends as the treated unit does. The remaining term concerns the counterfactual cycle $G^0_t=e_t-b'v_t$ of \eqref{eq:cycle0}, the deviation of the outcome the treated unit would have had without the policy from the counterfactual that its controls provide, and it measures how much of that deviation a policy surprise predicts. Assumption \ref{assu:A4} sets it to zero. The assumption does not require the deviation to vanish: the treated unit may still move away from its controls for reasons of its own, as long as the policy surprise does not predict those movements. In the example, the control states must provide a credible counterfactual for employment in Alaska, in the sense that the movements of Alaska away from that counterfactual which are not due to the dividend are unrelated to the dividend surprise.

\begin{thm}[Identification\label{thm:ident}]
Under Assumptions \ref{assu:A1} to \ref{assu:A5}, $\gamma_h=\theta_h=\sum_{j=0}^h\tau_j\psi_{h-j}$ for all $h\geq0$.
\end{thm}

The response of the cycle is therefore the effect of the policy on the treated unit: the impulse causal effects of the policy, combined with the path the policy follows after a surprise. When the effect of the policy varies over time, $\theta_h$ averages it with weights proportional to the surprise, and the average is an average treatment effect on the treated for a dummy policy and an average causal response for a continuous one (Supplementary Appendix \ref{sec:identification}).

Under the conditional exogeneity of Assumption \ref{assu:A4}, the CVAR therefore recovers the causal effect without bias, which a VAR without controls does not, as the following example shows.

\begin{example}
Consider a system with one control and one common trend, in which the control is $Y^0_t=F_t+v_t$, the treated unit without the policy is $Y^1_t(0)=F_t+e_t$, and the common trend follows $\Delta F_t=\eta_t$. The policy reacts to the common shock, $W_t=c\,\eta_t+\varepsilon_t$, and has only an impact effect, so that $Y^1_t=Y^1_t(0)+\tau_0W_t$, with all shocks independent, serially uncorrelated and of unit variance. The impact response of the VAR without controls is
\begin{equation*}
\gamma^V_0=\tau_0+\frac{c}{1+c^2}\qquad\text{(VAR)}.
\end{equation*}
If the policy has no effect, $\tau_0=0$, and reacts one for one to the common shock, $c=1$, the VAR reports an impact effect of one half. The cycle is $G_t=e_t-v_t+\tau_0W_t$, which does not contain the common shock, and the CVAR returns
\begin{equation*}
\gamma_0=\tau_0=0\qquad\text{(CVAR)}.
\end{equation*}
\end{example}

By ignoring the common shock, the VAR reports an effect where there is none, and the mistake has consequences for policy. A local government that responds to a downturn may conclude that its intervention worked, or that it failed, when the downturn and the recovery that followed came from a broader disruption in the sectors on which its economy depends. Comparing the treated unit with units of similar economic composition that went through the same disruption without the policy separates the two. The difference between the response of a VAR without controls and that of the CVAR then measures the gap between the conclusion that the policymaker reaches by observing the treated unit alone and the one reached by comparing it with units that did not adopt the policy.

The response of the cycle combines the effect of a policy surprise with that of the later values of the policy that follow it, and the same VAR separates the two, because it also estimates the path $\psi_h$ that the policy follows after a surprise.

\begin{thm}[Impulse causal effects and regime effects\label{thm:direct}]
Let Assumptions \ref{assu:A1} to \ref{assu:A5} hold.
\begin{enumerate}
\item[(a)] The impulse causal effects are identified recursively from the responses of the cycle and of the policy, $\tau_0=\gamma_0$ and $\tau_h=\gamma_h-\sum_{j=0}^{h-1}\tau_j\psi_{h-j}$ for $h\geq1$.
\item[(b)] If, in addition, Assumption \ref{assu:A3b} holds and $\E[W_{t+k}(0)\,u_t]=0$ for all $k\geq1$, then $\gamma_h=\theta^\ast_h$, the regime effect.
\end{enumerate}
\end{thm}

Part (a) of Theorem \ref{thm:direct} requires nothing about how the policy evolves after the surprise: whatever path the policy follows, the CVAR measures it, and removing it from the response of the cycle leaves the direct effects. Part (b) gives the condition under which no correction is needed: when the surprise does not predict the path the policy would have followed without it, the impulse response is already the regime effect.

\subsection{Bounds\label{sec:bounds}}

The controls never receive the policy, so under Assumption \ref{assu:A1} any response they show to its innovation comes from the shocks the policy reacted to, and that response is the part of the VAR estimate that a control removes. The controls are therefore useful beyond the construction of the cycle: they measure the bias of a VAR without controls and, as this subsection shows, they give a scale for the bias that the CVAR leaves when Assumption \ref{assu:A4} is in doubt. The gap between the two responses can be measured in a VAR that contains both the treated unit and the weighted controls. Let $\breve u_t$ be the innovation of the policy in a recursive VAR in $(R_t',W_t,Y^1_t,C_t)'$ with $R_t$ first and $W_t$ second, and for any variable $X_t$ of the system let $\gamma^X_h=\E[(X_{t+h}-X_{t-1})\breve u_t]/\E[\breve u_t^2]$ be its response.

\begin{prop}[Response of the controls\label{prop:placebo}]
\begin{enumerate}
\item[(a)] For every $h\geq0$, $\gamma^{Y^1}_h=\gamma^G_h+\gamma^C_h$.
\item[(b)] Under Assumptions \ref{assu:A1}, \ref{assu:A2} and \ref{assu:A5}, $\gamma^C_h=\lambda'\E[(F_{t+h}-F_{t-1})\breve u_t]/\E[\breve u_t^2]+b'\E[(v_{t+h}-v_{t-1})\breve u_t]/\E[\breve u_t^2]$. If the idiosyncratic components of the controls are uncorrelated with $\breve u_t$ at all leads and lags, $\gamma^C_h$ is the common-trend term of Theorem \ref{thm:var} computed with $\breve u_t$.
\item[(c)] Under Assumptions \ref{assu:A1} to \ref{assu:A5}, with $\breve u_t$ in place of $u_t$, $\gamma^G_h$ is the causal effect of a policy surprise.
\end{enumerate}
\end{prop}

A researcher who has estimated a VAR without controls can add the weighted controls to the system and read the size of the bias from their response, and a test that this response is zero at all horizons is a test of the absence of the common-trend bias. The response of the treated outcome splits exactly into the causal response of the cycle and the response of the controls.

The CVAR still rests on Assumption \ref{assu:A4} for the bias specific to the treated unit, and a reader may doubt it in a given application. The response of the controls also helps here, because it gives a natural scale for how large a violation would have to be to change a conclusion: instead of setting the bias specific to the treated unit to zero, one can bound it by a multiple of the bias from the common trends, which the controls measure. In the VAR of Proposition \ref{prop:placebo}, write the response of the cycle as $\gamma^G_h=\breve\theta_h+\breve B_h$, where $\breve\theta_h$ is the effect of a policy surprise and $\breve B_h=\E[G^0_{t+h}\breve u_t]/\E[\breve u_t^2]$ the bias specific to the treated unit, and let $\gamma^C_h$ be the response of the weighted controls, which measures the bias from the common trends.

\begin{assumptionc}[Bounded bias specific to the treated unit\label{assu:A4b}]
For a given $\lambda\geq0$, $|\breve B_h|\leq\lambda\,|\gamma^C_h|$ for every horizon $h$ of interest.
\end{assumptionc}

Assumption \ref{assu:A4b} replaces exogeneity with a limit on how much the policy surprise can be confounded by developments specific to the treated unit. The limit is set relative to a quantity that the data reveal, the response of the weighted controls, which measures how strongly the surprise is confounded by the common trends. The assumption allows the bias specific to the treated unit to reach at most a multiple $\lambda$ of that measured bias, in either direction. With $\lambda=1/2$, for example, suppose that employment in the control states falls after a dividend surprise, because of the national developments that come with the surprise. The developments specific to Alaska that come with the same surprise may then move employment in Alaska by at most half as much, up or down. The multiple concerns the confounding that comes with the surprise, and it says nothing about an effect of the policy on the controls, which Assumption \ref{assu:A1} excludes. Assumption \ref{assu:A4} is the case $\lambda=0$, and $\lambda=1$ allows the confounding specific to the treated unit to be as strong as the confounding from the common trends. The comparison follows the logic of \citet{AltonjiElderTaber2005} and \citet{Oster2019}, who bound the selection on unobserved variables by the selection on observed ones: here the observed confounding is the one that the common trends carry, which the controls reveal, and the unobserved confounding is the one specific to the treated unit.

\begin{prop}[Bounds\label{prop:bounds}]
Under Assumptions \ref{assu:A1}, \ref{assu:A2}, \ref{assu:A3}, \ref{assu:A5} and \ref{assu:A4b}, and if the idiosyncratic components of the controls are uncorrelated with $\breve u_t$ at all leads and lags:
\begin{enumerate}
\item[(a)] $\breve\theta_h\in\big[\gamma^G_h-\lambda|\gamma^C_h|,\ \gamma^G_h+\lambda|\gamma^C_h|\big]$, an interval that reduces to $\gamma^G_h$ when $\lambda=0$.
\item[(b)] The response of the VAR without controls, $\gamma^{Y^1}_h=\gamma^G_h+\gamma^C_h$, belongs to the interval only if $\lambda\geq1$: attributing to the policy the response of a VAR without controls requires a bias specific to the treated unit as large as the bias from the common trends.
\end{enumerate}
\end{prop}

The bounds are estimated by the corresponding responses of the VAR of Proposition \ref{prop:placebo}, and bands for their end points follow from bootstrapping that VAR. The breakdown value of $\lambda$, at which a conclusion stops holding, summarises how strong the violation of Assumption \ref{assu:A4} would have to be to overturn it, as in the sensitivity analysis of \citet{RambachanRoth2023} for parallel trends.

\subsection{Estimation and inference\label{sec:estimation}}

The theorems take the weights on the controls as given. In practice they are unknown, so the response of the CVAR is computed from a cycle built with estimated weights, and inference has to account for the fact that a variable of the VAR is itself estimated. The problem is familiar in econometrics, where it arises with generated regressors \citep{Pagan1984} and with two-step estimators in general \citep{MurphyTopel1985,NeweyMcFadden1994}: the sampling error of the first step enters the distribution of the second unless it vanishes faster than the error of the second step. In the CVAR the first step has a feature of its own, because the weights come from a cointegrating regression, in which the part that multiplies the common trends is estimated faster than the rest. I estimate the weights and the response in four steps: the first two build the counterfactual from the controls, and the last two estimate the VAR on the resulting cycle.

\textit{Step 1.} I estimate how many common trends the controls share. When there are more controls than trends, the controls are cointegrated among themselves, and the number of their cointegrating relations, $K-\hat r$, has to be estimated, as in a cointegrated VAR \citep{Johansen1991}. The difference lies in the use of the estimate. In a cointegrated VAR the cointegrating relations are the object of interest, while in the CVAR they tell the next step which combinations of the controls carry the trends and which are stationary and serve only to absorb idiosyncratic variation. I estimate $K-\hat r$ with an information criterion on the eigenvalues of the reduced-rank regression of \citet{Johansen1991}, with a penalty that grows with the logarithm of the sample size \citep{AznarSalvador2002,ChengPhillips2009}, collect the estimated cointegrating vectors in $\hat H$, and let $\hat D_1$ span the orthogonal complement of $\hat H$, with $\hat D_1=I_K$ when $\hat r=K$. The combinations $\hat X_{1t}=\hat D_1'Y^0_t$ carry the common trends of the controls. The usual estimation of a cointegrating relation treats the regressors as having independent trends and has no such step.

\textit{Step 2.} I estimate the weights $(a,b)$ by dynamic least squares \citep{Saikkonen1991,StockWatson1993} in the regression
\begin{equation}
Y^1_t=a+b'Y^0_t+\sum_{j=0}^{q}\kappa_j'\begin{pmatrix}R_{t-j}\\W_{t-j}\end{pmatrix}+\sum_{j=-L}^{L}\delta_j'\Delta\hat X_{1,t+j}+\varepsilon_t,
\label{eq:dols}
\end{equation}
optionally with the constraint $b\geq0$. Two features of \eqref{eq:dols} depart from the standard regression. The policy terms are included because, by \eqref{eq:cycle}, the error of a regression of $Y^1_t$ on $Y^0_t$ alone would contain the effect of the policy, which is correlated with the innovations of the trends whenever the policy reacts to them. The leads and lags are taken only in the directions of the trends, because leads and lags of the stationary combinations of the controls would span those combinations as $L$ grows and leave their weights unidentified. The levels enter as $b'Y^0_t$, so the fit does not depend on the choice of $\hat D_1$.

\textit{Step 3.} With the weights in hand, I compute the cycle $\hat G_t=Y^1_t-\hat a-\hat b'Y^0_t$. The estimated policy terms are not subtracted, since the cycle must contain the effect of the policy.

\textit{Step 4.} I estimate a VAR($p$) with a constant in $\hat Z_t=(R_t',W_t,\hat G_t)'$, and compute the response of $\hat G_t$ to an innovation in $W_t$ in the recursive identification with $R_t$ first, normalised to a unit impact on the policy.

I call the weights that the estimator selects in the limit the pseudo-true weights $b^\ast$, and $G^\ast_t=Y^1_t-b^{\ast\prime}Y^0_t$ the cycle they produce, for which Assumption \ref{assu:A4} is required.\footnote{Formally, write the controls in the rotated coordinates $D'Y^0_t=(X_{1t}',S_t')'$ with $D=(D_1,H)$, so that $b'Y^0_t=c_1'X_{1t}+c_2'S_t$ with $c=D^{-1}b$; $c_1$ multiplies the trends and $c_2$ the stationary combinations. The pair $(a^\ast,b^\ast)$ is the limit, as $L\to\infty$, of the population minimisers of the squared error of \eqref{eq:dols} over $b\in\mathcal{B}$, with the constraint when it is imposed. It exists and is unique when the spectral density of $(\Delta X_{1t}',S_t',R_t',W_t)'$ is bounded away from singularity.}

\begin{thm}[Estimated weights\label{thm:est}]
Let Assumptions \ref{assu:A1} to \ref{assu:A5} hold with $b=b^\ast$, let the spectral density of $(\Delta X_{1t}',S_t',R_t',W_t)'$ be bounded away from singularity, let $L\to\infty$ with $L^3/T\to0$, let the regularity conditions of \citet{Saikkonen1991} hold for \eqref{eq:dols},\footnote{The vector of the changes of the trends and of the error of \eqref{eq:dols} with all leads and lags, $(\Delta X_{1t}',\varepsilon^\infty_t)'$, is a stationary linear process with absolutely summable coefficients and independent innovations with finite fourth moments, and its spectral density at frequency zero is positive definite. The coefficients $\delta_j$ of the projection of $\varepsilon^\infty_t$ on all leads and lags of $\Delta X_{1t}$ satisfy $\sqrt{T}\sum_{|j|>L}\|\delta_j\|\to0$, which makes the omitted leads and lags negligible at the rate $T^{-1/2}$ that the coefficients on the stationary regressors require.} let $(R_t',W_t,G^\ast_t,\Delta X_{1t}',S_t')'$ be strictly stationary and strong mixing with coefficients $\alpha(m)$ such that $\sum_{m\geq1}\alpha(m)^{\epsilon/(2+\epsilon)}<\infty$ and have finite moments of order $4+2\epsilon$ for some $\epsilon>0$, let $\mathbb{P}(\hat r=r)\to1$,\footnote{The information criterion of the first step satisfies this condition \citep{AznarSalvador2002,ChengPhillips2009}. The trace test of \citet{Johansen1991} with a fixed level selects the true rank with probability approaching one minus the level, and with it the conclusions of the theorem hold on an event of that probability.} and let $\hat c=D^{-1}\hat b$ and $c^\ast=D^{-1}b^\ast$.
\begin{enumerate}
\item[(a)] $\hat c_1-c_1^\ast=O_p(T^{-1})$, $\hat c_2-c_2^\ast=O_p(T^{-1/2})$ and $\hat a-a^\ast=O_p(T^{-1/2})$. Without the constraint, or when $b^\ast$ is in the interior of the constraint set, $\sqrt{T}(\hat c_2-c_2^\ast)$ is asymptotically normal.
\item[(b)] If $Z_t$ follows a stationary VAR($p$) with $G^\ast_t$ as the cycle, let $\tilde\gamma_h$ be the estimator of the response with $(a^\ast,b^\ast)$ known. Then, for every fixed $h$,
\begin{equation}
\sqrt{T}(\hat\gamma_h-\gamma_h)=\sqrt{T}(\tilde\gamma_h-\gamma_h)+J_h'\sqrt{T}(\hat c_2-c_2^\ast)+o_p(1),
\label{eq:firststep}
\end{equation}
where $J_h$ is the derivative of the population response with respect to the weights on $S_t$. Without the constraint, or when $c^\ast$ is in the interior of the constraint set, the two terms on the right are jointly asymptotically normal.
\item[(c)] If $r=K$, the second term in \eqref{eq:firststep} is absent, with or without the constraint.
\end{enumerate}
\end{thm}

The error in the weights on the trends is of order $T^{-1}$ and disappears from the distribution of the response, whereas the weights on the stationary combinations of the controls, which absorb idiosyncratic variation of the treated unit as synthetic control weights do, are estimated at the same rate as the coefficients of the VAR, so their error enters the distribution of the response.

The error of these synthetic control weights enters the response only indirectly, through the cycle on which the VAR is estimated and the policy innovation it implies. Proposition \ref{prop:J} in Supplementary Appendix \ref{sec:supp-inf} gives its coefficient, which is zero when the policy surprise $u_t$ is orthogonal to the stationary movements of the controls $S_{t+j}$ at every lead and lag, as it is when the controls never receive the policy and their idiosyncratic shocks are unrelated to the policy of the treated unit, and when the policy does not react to the past of the cycle, so that a change in the weights leaves the policy innovation unchanged.

When the conditions fail, inference must account for the first step. I compute bands with a residual bootstrap that regenerates the controls together with the policies and the cycle, and re-estimates the weights in every draw. Let $\hat H$ span the estimated cointegrating vectors of the controls and $\hat D_1$ its orthogonal complement, as in the first estimation step, and let $\hat X_{1t}=\hat D_1'Y^0_t$ and $\hat S_t=\hat H'Y^0_t$. The bootstrap proceeds as follows.
\begin{enumerate}
\item[(i)] Estimate a VAR($p$) with a constant in $\hat Q_t=(\Delta\hat X_{1t}',\hat S_t',R_t',W_t,\hat G_t)'$ and collect the centred residual vectors.
\item[(ii)] In each draw, resample the residual vectors with replacement and rebuild $(\Delta X^\ast_{1t},S^\ast_t,R^\ast_t,W^\ast_t,G^\ast_t)$ recursively from the observed initial values.
\item[(iii)] Rebuild the controls, cumulating $\Delta X^\ast_{1t}$ from the observed initial value and rotating back, $Y^{0\ast}_t=(\hat D_1,\hat H)'^{-1}(X^{\ast\prime}_{1t},S^{\ast\prime}_t)'$, and the treated outcome, $Y^{1\ast}_t=\hat a+\hat b'Y^{0\ast}_t+G^\ast_t$.
\item[(iv)] Re-estimate the weights by \eqref{eq:dols} on the bootstrap sample, with the rank and $\hat D_1$ held at their estimates, compute the bootstrap cycle, estimate the VAR($p$) in $(R^\ast_t,W^\ast_t,\hat G^\ast_t)$ and compute the response.
\item[(v)] Simulate one long path of the VAR of step (i), many times longer than the sample, apply steps (iii) and (iv) to it, and call the resulting response $\gamma^\ast_h$; shift every draw by $\hat\gamma_h-\gamma^\ast_h$.
\end{enumerate}
The bands are percentile bands of the shifted draws.

The bootstrap of the CVAR differs from that of a standard or a cointegrated VAR in two respects. In a VAR estimated by least squares, or a cointegrated VAR estimated by maximum likelihood, the parameters of the bootstrap world are the estimates themselves, so the bootstrap distribution is centred at the estimate by construction. The CVAR adds a first step: the weights that remove the common trends are a projection of the treated unit on the controls, and they determine the cycle on which the VAR is estimated. The error in the weights therefore enters the residuals of the cycle, and through them the response. A bootstrap that holds the controls at their observed values cannot reproduce this error, because it regenerates a cycle that is independent of the observed controls.\footnote{By the argument of Proposition \ref{prop:J}, the first-step term of such a bootstrap vanishes asymptotically, so the bootstrap is valid only when $J_h=0$.} Steps (ii) and (iii) regenerate the controls together with the policies and the cycle, which reproduces the dependence between them that Proposition \ref{prop:J} makes explicit. The second respect is the centring. In the bootstrap world generated by the joint VAR, the weights to which the projection converges, and the response they produce, are functions of the autocovariances of the fitted joint VAR, which differ from those of the sample, so the bootstrap distribution is centred at a value that differs from the estimate. The problem is the one that \citet{HallHorowitz1996} solve for overidentified GMM estimators, whose moment conditions do not hold in the bootstrap world \citep{Horowitz2001}, and it is close to the one that \citet{Kilian1998} addresses when he estimates the bias of impulse responses in the bootstrap world. Step (v) applies their recentring principle: it computes the response of the bootstrap world by simulation, as indirect inference computes its binding function \citep{GourierouxMonfortRenault1993}, and centres the draws at it. When the weights are known, a bootstrap of the VAR in $Z_t$ alone suffices.

With the weights known, the residual bootstrap of the VAR in $Z_t$ gives percentile bands with asymptotic coverage $1-\alpha$ by the standard argument for autoregressions (Theorem \ref{thm:resid} in Supplementary Appendix \ref{sec:supp-inf}).

The bootstrap, like the estimator, depends on the lag length of the VAR in the fourth step. The counterfactual cycle sums the idiosyncratic components of the treated unit and of the controls, so $Z_t$ is a VARMA whose autoregressive representation has infinite order and slowly declining coefficients, and the relevant theory lets the lag length grow with the sample, $p_T\to\infty$ with $p_T^3/T\to0$ \citep{LewisReinsel1985,InoueKilian2002}. Information criteria select too few lags in designs of this kind (Supplementary Appendix \ref{sec:montecarlo}), so I use one year of monthly lags, and when the sample is too short for long lags, local projections of the cycle on the policy innovation estimate the same response \citep{PlagborgMollerWolf2021,LiPlagborgMollerWolf2024}.

With the weights estimated and a lag length that grows with the sample, the following result covers both the estimator and the bootstrap of steps (i) to (iv).

\begin{thm}[Estimated weights, long autoregressions and the bootstrap\label{thm:joint}]
Let the conditions of Theorem \ref{thm:est}, except the finite order of the VAR in part (b), hold, with $c^\ast$ in the interior of the constraint set when the constraint is imposed.\footnote{When a weight is on the boundary, the limit in part (a) is the projection of a normal vector on a cone, and the bootstrap is inconsistent \citep{Andrews2000}.} Let $Q_t=(\Delta X_{1t}',S_t',R_t',W_t,G^\ast_t)'$ have the representation $Q_t=\mu+\sum_{l\geq1}\Phi_lQ_{t-l}+\epsilon_t$, with $\det(I-\sum_{l\geq1}\Phi_lz^l)\neq0$ for $|z|\leq1$, $\sum_{l\geq1}l\|\Phi_l\|<\infty$, and i.i.d. innovations $\epsilon_t$ with finite fourth moments and positive definite covariance, and let $Z_t$ have a representation with the same properties. Let the lag length of the VARs in $\hat Z_t$ and in $\hat Q_t$ be $p=p_T$, with $p_T\to\infty$, $p_T^3/T\to0$ and $\sqrt{T}\sum_{l>p_T}\|\Phi_l\|\to0$ and the same for the coefficients of $Z_t$, and let the number of leads and lags $L$ and of lags of the policies $q$ in \eqref{eq:dols} be at most $p_T$. Then, for every fixed $h$:
\begin{enumerate}
\item[(a)] the expansion \eqref{eq:firststep} of Theorem \ref{thm:est}(b) holds, with $\gamma_h$ the response of the VAR($\infty$) representation of $Z_t$ and $J_h$ its derivative with respect to the weights on $S_t$, and the two terms on its right are jointly asymptotically normal;
\item[(b)] the distribution of $\sqrt{T}(\hat\gamma^\ast_h-\gamma^\ast_h)$ under the bootstrap of steps (i) to (iv), where $\gamma^\ast_h$ is the response in the bootstrap world, converges in probability to the limit of the distribution of $\sqrt{T}(\hat\gamma_h-\gamma_h)$, and the percentile band of the draws shifted in step (v) has asymptotic coverage $1-\alpha$ as the length of the long path grows with $T$.
\end{enumerate}
\end{thm}

Part (b) holds whether or not $J_h$ is zero,\footnote{Theorems \ref{thm:resid} and \ref{thm:joint} assume i.i.d. innovations. Under conditional heteroskedasticity of unknown form, resampling residual vectors one at a time does not reproduce the fourth moments that enter the variance of the innovation covariance, and the residual-based moving-block bootstrap of \citet{bruggemann2016inference} gives valid inference for the responses of a VAR at every horizon, including impact; it can replace the resampling of step (ii). The wild bootstrap reproduces the variance of the slope coefficients \citep{goncalveskilian2004}, and Supplementary Appendix \ref{sec:supp-inf} shows that it does not reproduce that of the innovation covariance.} so the bands account for the estimation of the weights in every case covered by Proposition \ref{prop:J}. When $r=K$ there are no stationary combinations of the controls, $S_t$ drops out of $Q_t$, and the first-step term is absent by Theorem \ref{thm:est}(c). In the simulations of Supplementary Appendix \ref{sec:montecarlo}, the joint bootstrap gives bands with coverage close to that obtained with the weights known at every horizon up to a year. The simulations also show that a wild bootstrap with Rademacher weights is not a valid alternative at impact (Supplementary Appendix \ref{sec:supp-inf}).\footnote{Theorem \ref{thm:wild} in Supplementary Appendix \ref{sec:supp-inf} shows that, with Rademacher weights, the bootstrap covariance of the residuals differs from the estimate by a term of smaller order than its sampling error, so the bands at impact are too narrow.}

\section{Testing the design\label{sec:test}}

The CVAR rests on Assumption \ref{assu:A2}: the treated unit shares the stochastic trends of the controls, so that the counterfactual cycle is stationary. When the assumption fails, the treated unit carries a trend of its own, the cycle has a unit root, and the responses of Section \ref{sec:theorems} are spurious. Testing the assumption is a test of the null hypothesis of no cointegration between the treated unit and its controls. The standard tests for this null estimate the cointegrating vector first and test the residual for a unit root \citep{EngleGranger1987,PhillipsOuliaris1990}, possibly after GLS detrending \citep{PerronRodriguez2016}, and the stationarity tests reverse the null \citep{Shin1994,HarveyThiele2021}. Three features of the CVAR make their tables inapplicable: the weights are estimated together with the policy terms, they may be constrained, and the controls may be cointegrated among themselves. The simulations of Supplementary Appendix \ref{sec:montecarlo} show that the residual-based tests then reject a true null far more often than their nominal level. The power of tests of this null has been studied by \citet{Pesavento2004}, \citet{ElliottJanssonPesavento2005} and \citet{ElliottPesavento2009}, who show that no uniformly most powerful test exists and derive point-optimal tests and power bounds. I adapt the point-optimal approach to the CVAR.

Write the treated unit as
\begin{equation}
Y^1_t=\mu+b'Y^0_t+\delta(L)'P_t+u_t,\qquad u_t=\rho_Tu_{t-1}+v_t,\qquad \rho_T=1+c/T,
\label{eq:testmodel}
\end{equation}
where $P_t$ collects the current value and $q$ lags of the policies and $v_t$ is a stationary linear process with long-run variance $\omega^2$. The null hypothesis $c=0$ is that the treated unit has a stochastic trend of its own, and the CVAR is not valid; the alternatives $c<0$ are that the cycle is stationary. For $\alpha\in(0,1]$ let $S(\alpha)$ be the sum of squared residuals of the regression of $Y^1_t-\alpha Y^1_{t-1}$ on the quasi-differenced constant, controls and policy terms, with the first observation in levels, and let $\hat\omega^2$ be the autoregressive spectral estimator of $\omega^2$ from the augmented Dickey-Fuller regression on the residuals of the regression with $\alpha=\bar\alpha$, with the lag order chosen by the modified information criterion of \citet{NgPerron2001}. For a fixed $\bar c<0$ and $\bar\alpha=1+\bar c/T$, the test rejects for small values of
\begin{equation}
P_T=\frac{S(\bar\alpha)-\bar\alpha\,S(1)}{\hat\omega^2}.
\label{eq:PT}
\end{equation}
The regression is unconstrained: the restriction $b\geq0$ belongs to the estimator and plays no role in the test. The weights are concentrated out of the quasi-differenced regression under both hypotheses, which is the flat-weight limit of the family of tests of \citet{ElliottPesavento2009}. The difference with the residual-based tests is that the controls enter the regression quasi-differenced at the same $\alpha$ as the treated unit, so that their coefficient is identified mainly from the co-movement of the changes, under the null and the alternative alike.

The properties of the test rest on Assumption \ref{assu:T} in Supplementary Appendix \ref{sec:supp-testboot}. It requires the policies to be stationary under both hypotheses, the shocks specific to the treated unit to be unrelated to the changes of the controls and of the policies at all leads and lags, which excludes a reaction of the policy to the past of the component specific to the treated unit, and a small initial condition.

\begin{thm}[Asymptotic equivalence with the known-weights test\label{thm:po}]
Let Assumption \ref{assu:T} hold and let $P^u_T$ be the statistic \eqref{eq:PT} computed from $u_t$ itself, with only the constant as regressor. Under the null and under local alternatives $c$ fixed, $P_T=P^u_T+o_p(1)$. Consequently:
\begin{enumerate}
\item[(a)] $P_T\Rightarrow\bar c^2\int_0^1W_c(s)^2ds-\bar c\,W_c(1)^2$, where $W_c$ is an Ornstein-Uhlenbeck process with parameter $c$ and $W_0$ a standard Brownian motion; under the null the limit does not depend on the number of controls $K$, on the number of their trends $r$, or on the policy terms, and the critical values are those of \citet{ERS1996} for the case with a constant.\footnote{With $\bar c=-7$, the choice for which the envelope has power one half at the 5 percent level \citep[p.~821]{ERS1996}, the asymptotic critical values at 1, 2.5, 5 and 10 percent are 1.99, 2.55, 3.26 and 4.48 \citep[Table I]{ERS1996}.}
\item[(b)] The asymptotic local power of the test equals that of the point-optimal test computed from the counterfactual cycle with known weights and policy coefficients. At $c=\bar c$ it equals the asymptotic Gaussian power envelope, which bounds the power of any test of a unit root in the cycle \citep{ERS1996}; since the test does not use the weights, it attains at $\bar c$ the envelope of the tests that do not know them.
\end{enumerate}
\end{thm}

Part (a) says that estimating the weights inside the test leaves its null distribution unchanged in large samples, which is what distinguishes it from the residual-based tests of the null of no cointegration. Those tests estimate the weights by a regression in levels, and under the null the regression is spurious: least squares chooses the weights that make the residual as small as possible, which also makes it look closer to stationarity than the true deviation of the treated unit from its controls, and the more so the more controls the regression uses. The critical values of the residual-based tests therefore depend on the number of controls in the regression, and a separate table is needed for each number of controls \citep{PhillipsOuliaris1990,PerronRodriguez2016}. In the point-optimal test the controls enter quasi-differenced at $\bar\alpha$, close to one, so their coefficients are identified mainly from the co-movement of the changes, where there is no spurious relation to fit, and the estimated weights cannot make a random walk look stationary. The statistic then behaves as if it were computed from the true deviation $u_t$, and its critical values are those of the univariate test of \citet{ERS1996} with a constant, whatever the number of controls, of their trends or of the policy terms. Part (b) means that, under Assumption \ref{assu:T}, not knowing the weights costs no power in large samples: the test is as powerful as if the counterfactual cycle were observed. In finite samples the estimation of the weights and of the policy coefficients shifts the null distribution slightly, and the long-run variance of a persistent cycle is estimated with error. I therefore compute the p-value with a sieve bootstrap under the null, which rebuilds the treated unit by cumulating its fitted changes and resampled autoregressive errors, holds the controls and the policy terms at their observed values, and has asymptotic size equal to its nominal level (Proposition \ref{prop:sieve} in Supplementary Appendix \ref{sec:supp-testboot}).



The test decides whether the more efficient estimator in levels can be used. Estimated on the changes of the cycle, the CVAR identifies the effect whether or not the treated unit has a trend of its own, so a non-rejection costs efficiency and leaves the identification unchanged. A rejection supports Assumption \ref{assu:A2}, and the CVAR can then be estimated in levels with the weights of \eqref{eq:dols}. The test is consistent, so when the cycle is stationary it selects the estimator in levels with probability approaching one.\footnote{When the cycle has a trend of its own, the test selects the estimator in levels with probability equal to its size, so the size of the test also sets the protection of the estimator in levels against a stochastic trend in the cycle.} The power of any test of this null is low against a cycle that returns to its mean slowly relative to the length of the sample, because such a cycle is hard to tell apart from a random walk. Supplementary Appendix \ref{sec:screen} applies the test to the states of the United States, each compared with its neighbouring states, and reports its power in panels built like those data.

\section{A treated unit with a trend of its own\label{sec:ext}}

Assumption \ref{assu:A2} asks the weighted controls to reproduce every stochastic trend of the treated unit. When the test of Section \ref{sec:test} does not reject, the treated unit may carry, alongside the trends it shares with the controls, a slow-moving component of its own that no combination of the controls reproduces. In the application, the national cycle and the price of oil move employment in every state, and Alaska can also follow a path of its own over four decades, through the rise and decline of oil production on the North Slope. The controls still remove the shared trends, and the CVAR can be estimated on the changes of the cycle, which is the version I use in the applications.

\begin{assumption}[Common trends and a trend of the treated unit\label{assu:A2d}]
There is an $r$-dimensional process $F_t$, integrated of order one with stationary first difference, a $K\times r$ matrix $\Lambda$ of rank $r\leq K$ and a known vector $b\in\mathbb{R}^K$ such that
\begin{equation*}
Y^0_t=\Lambda F_t+v_t,\qquad Y^1_t(0)=\lambda'F_t+e_t,\qquad \lambda=\Lambda'b,
\end{equation*}
where $v_t$ is stationary and $e_t$ is integrated of order at most one, with stationary first difference.
\end{assumption}

Assumption \ref{assu:A2d} keeps the condition that the weighted controls share the trends of the treated unit and drops the stationarity of $e_t$. The counterfactual cycle $G^0_t=e_t-b'v_t$ can then be integrated of order one, and its change $\Delta G^0_t$ is stationary. The weights are set by the researcher or estimated. When $e_t$ has a stochastic trend, a regression of the treated outcome on the controls in levels has no well-defined limit, and a regression in changes estimates the projection of $\Delta Y^1_t$ on $\Delta Y^0_t$, in which the changes of $v_t$ act as measurement error in the regressors and move the weights away from those that remove the trends. In the application to Alaska I estimate the weights from the changes. Under Assumptions \ref{assu:A1}, \ref{assu:A2d} and \ref{assu:A3}, the decomposition \eqref{eq:cycle} holds unchanged, because it uses only $\lambda=\Lambda'b$.

Let $Z^\Delta_t=(R_t',W_t,\Delta G_t)'$, let $u_t$ be the innovation of the policy in the recursive VAR in $Z^\Delta_t$, defined as in \eqref{eq:innovation} with the past of $\Delta G_t$ in place of that of $G_t$, let $\theta_h$ and $\psi_k$ be defined as in \eqref{eq:estimand} and \eqref{eq:decomp} with this innovation, and let
\begin{equation}
\Gamma_h=\frac{\E[(G_{t+h}-G_{t-1})u_t]}{\E[u_t^2]}
\label{eq:cumcycle}
\end{equation}
be the response of the cycle, cumulated from the month before the surprise.

\begin{assumption}[Exogeneity with respect to the changes of the counterfactual cycle\label{assu:A4d}]
(i) $(R_t',W_t,\Delta G^0_t,\Delta F_t',\Delta v_t')'$ is covariance stationary with finite fourth moments, and $\E[u_t^2]>0$. (ii) $\E[(G^0_{t+h}-G^0_{t-1})u_t]=0$ for all $h\geq0$.
\end{assumption}

Part (i) plays the role of Assumption \ref{assu:A5} for the changes of the cycle. Part (ii) plays the role of Assumption \ref{assu:A4}: it allows the policy to react to the past changes of the cycle, which the VAR conditions on, and requires that the surprise does not predict how the component specific to the treated unit moves from the month before the surprise onward. In the application, it requires that a surprise in the dividend carries no news about the path that employment specific to Alaska would have followed without it.

\begin{thm}[Response of the cumulated cycle\label{thm:diff}]
Let Assumptions \ref{assu:A1}, \ref{assu:A2d}, \ref{assu:A3} and \ref{assu:A4d}(i) hold.
\begin{enumerate}
\item[(a)] For every $h\geq0$,
\begin{equation*}
\Gamma_h=\theta_h+\frac{\E[(G^0_{t+h}-G^0_{t-1})u_t]}{\E[u_t^2]}.
\end{equation*}
If $Z^\Delta_t$ follows a VAR, possibly of infinite order, $\Gamma_h=\sum_{j=0}^{h}\gamma^\Delta_j$, where $\gamma^\Delta_j$ is the response of $\Delta G_t$ at horizon $j$ to a unit innovation in $W_t$ in the recursive identification with $R_t$ first, $W_t$ second and $\Delta G_t$ last.
\item[(b)] If Assumption \ref{assu:A4d}(ii) also holds, $\Gamma_h=\theta_h$ for all $h\geq0$, and Theorem \ref{thm:direct} holds with $\Gamma_h$ in place of $\gamma_h$.
\item[(c)] If the imposed weights $b$ do not satisfy $\lambda=\Lambda'b$, then for every $h\geq0$
\begin{equation*}
\Gamma_h=\theta_h+(\lambda-\Lambda'b)'\frac{\E[(F_{t+h}-F_{t-1})u_t]}{\E[u_t^2]}+\frac{\E[(G^0_{t+h}-G^0_{t-1})u_t]}{\E[u_t^2]},
\end{equation*}
with $G^0_t=e_t-b'v_t$.
\end{enumerate}
\end{thm}

Theorem \ref{thm:diff} is the counterpart of Theorems \ref{thm:cycle} and \ref{thm:ident} for a cycle with a stochastic trend. The controls remove the common trends exactly as before when the imposed weights are the right ones, and part (c) gives the bias when they are not: the common-trend term of \eqref{eq:bias}, with the difference between the loadings of the treated unit and those of the weighted controls in place of the loadings of the treated unit. The VAR is estimated on the changes of the cycle, and the assumption that closes the identification concerns the changes of the counterfactual cycle. The response of a VAR without controls in Theorem \ref{thm:var} is written in cumulated changes, so the comparison between the two estimators carries over: the response of the VAR without controls in $(R_t',W_t,\Delta Y^1_t)'$ equals $\Gamma_h$ plus the common-trend term of \eqref{eq:bias}, up to the difference between the two innovations. Proposition \ref{prop:placebo} and the bounds of Proposition \ref{prop:bounds} are also stated in cumulated changes, and they hold with the VAR in $(R_t',W_t,\Delta Y^1_t,\Delta C_t)'$.

With known weights, the cycle $G_t=Y^1_t-b'Y^0_t$ is observed and there is no first step. The vector $Z^\Delta_t$ is stationary, so Theorem \ref{thm:resid} applies to the responses $\gamma^\Delta_j$ and, because $\Gamma_h$ is a finite sum of them, to $\Gamma_h$ at every fixed horizon: the residual bootstrap of the VAR in $Z^\Delta_t$, with the responses cumulated in each draw, gives percentile bands with asymptotic coverage $1-\alpha$. The change of the cycle is far less persistent than the cycle in levels, so shorter orders suffice: in the applications I use twelve lags or the smallest order at which the residuals show no autocorrelation, with longer orders as a robustness check.\footnote{The bands treat the selected order as fixed, as is standard in applied work with VARs.}

When the weights are estimated, the regression runs on the changes of the series, where the idiosyncratic movements of the controls are as large as their common movements, and the estimated weights inherit them.

\begin{prop}[Weights estimated in changes\label{prop:wchanges}]
Let Assumptions \ref{assu:A2d} and \ref{assu:A4d}(i) hold, let the policies affect $\Delta Y^1_t$ only through the current and lagged values that the regression projects out, and, after that projection, let $\Delta F_t$, $\Delta v_t$ and $\Delta e_t$ be mutually uncorrelated, with $\operatorname{Var}(\Delta F_t)=\Sigma_F$ and $\operatorname{Var}(\Delta v_t)=D$ positive definite. The least-squares weights of the regression of $\Delta Y^1_t$ on $\Delta Y^0_t$ converge in probability to
\begin{equation*}
w^\ast=(\Lambda\Sigma_F\Lambda'+D)^{-1}\Lambda\Sigma_F\lambda,
\end{equation*}
and the response of the cumulated cycle built with $w^\ast$ is given by part (c) of Theorem \ref{thm:diff} with $b=w^\ast$ and
\begin{equation*}
\lambda-\Lambda'w^\ast=\big(I_r+\Lambda'D^{-1}\Lambda\,\Sigma_F\big)^{-1}\lambda .
\end{equation*}
\end{prop}

The weighted controls reproduce the part of the common movements of the treated unit that their changes reveal above their own noise, and the rest stays in the cycle. With a single common trend, the weights are proportional to $D^{-1}\Lambda$, so their composition favours the controls that follow the trend most closely relative to their noise, and their scale is shrunk by the factor $s/(1+s)$, with $s=\Sigma_F\Lambda'D^{-1}\Lambda$; the loading left in the cycle is $\lambda/(1+s)$, and the term it adds to the response equals $1/s$ times the response of the weighted controls in Proposition \ref{prop:placebo}, with the same sign. The term vanishes as $\Lambda'D^{-1}\Lambda$ grows, when the controls are many and their idiosyncratic movements independent or when those movements are small. It is the counterpart in changes of the bias of synthetic control with an imperfect fit before the treatment \citep{FermanPinto2021}, which disappears as the number of controls grows \citep{Ferman2021}. Consistent estimation of the weights in changes with a small set of controls is an open problem, and I leave it to future work. The size of the term can be measured. With a large panel of untreated units, the principal components of their changes estimate the common movements \citep{Bai2003}, and the ratio of the loading of the weighted controls to that of the treated unit on each component gives the share of that movement the weights reproduce. Equal weights avoid the estimation noise, and the same ratio checks whether they reproduce the loadings of the treated unit.

\section{The Alaska Permanent Fund Dividend\label{sec:alaska}}

I apply the CVAR to three policies. The first is the dividend of the Alaska Permanent Fund, whose size changes from year to year with the income of the fund. The second, extended unemployment benefits, switches on in recessions, and it is the standard example of a policy that reacts to the cycle it is meant to affect. The third, the Low Carbon Fuel Standard of California, is a cost that changes every month and that fuel suppliers can pass on to consumers. In all three I estimate the CVAR in changes; the test of Section \ref{sec:test} does not reject, so the data do not support the estimator in levels. In the first I estimate the weights on the controls, and in the other two I give them equal weights.

Since 1982 the State of Alaska has paid every resident a yearly dividend from the Alaska Permanent Fund, a sovereign wealth fund built from the royalties of oil production. A public formula sets the total distribution at about a tenth of the income that the fund earned over the last five fiscal years, the total is divided equally among the eligible residents, and the payment arrives in a single transfer in October.\footnote{The first dividend was paid in June 1982. From 2016 the legislature has set the amount below the formula. In 2008 a one-time energy rebate of \$1,200 was paid together with the dividend.} The dividend per person has ranged from a few hundred to more than two thousand dollars.

The dividend is the largest permanent and universal cash transfer in the United States, and it has become the reference case for the effects of a universal basic income. Its effect on spending is well documented. Households raise their spending on nondurables in the quarter in which the payment arrives and do not respond to news about its size before then \citep{Kueng2018}, a response that an earlier study based on survey income had not detected \citep{Hsieh2003}. Its effect on the labour market is less settled, because two channels push in opposite directions. A cash transfer reduces the labour supply of the household that receives it \citep{Cesarini2017}, while the spending it finances raises the demand for labour in the sectors that sell to local customers. Comparing Alaska with a synthetic control built from other states, \citet{JonesMarinescu2022} find that the introduction of the dividend left the employment rate unchanged and raised part-time work, and they attribute the absence of an employment loss to the demand channel offsetting the supply channel.

Two features of the dividend make the comparison of Alaska with other states informative and a time-series analysis of Alaska alone misleading. The first is that only Alaska pays it, so the other states go through the same national cycle without it. The second is that its size follows the returns of a fund invested in a diversified portfolio of financial assets, which move with the stock market and the national business cycle. The largest dividends came after years of high returns, and those years preceded the recessions of 2001 and 2008. A VAR on Alaska data alone attributes to the dividend the national downturns that followed (Theorem \ref{thm:var}), while the CVAR compares Alaska with states that went through the same cycle without the dividend. The synthetic control of \citet{JonesMarinescu2022} rests on a single event, the introduction of the dividend, with five years of annual data before it, and estimates the average effect of the programme over the following decades. The CVAR uses the variation in the size of more than thirty dividends, estimates the weights on the controls over the whole sample, payment years included, and traces the response of the labour market month by month.

I use monthly data from January 1976 to December 2019, which end before the pandemic. The outcomes are the employment-population ratio, the unemployment rate and the participation rate from the Local Area Unemployment Statistics, and the logarithm of payroll employment in total, in leisure and hospitality and in retail trade from the Current Employment Statistics, which start in January 1990.\footnote{Seasonally adjusted series from the Bureau of Labor Statistics: \texttt{LASST020000000000007}, \texttt{LASST020000000000003} and \texttt{LASST020000000000008} for Alaska, \texttt{SMS02000000000000001}, \texttt{SMS02000007000000001} and \texttt{SMS02000004200000001}, and the analogous codes for the controls.} The policy $W_t$ is the dividend paid in the month, in thousands of 2024 dollars per person, and zero in the other months.\footnote{Amounts from the Permanent Fund Dividend Division, Summary of Applications and Payments, including the 2008 energy rebate, deflated by the consumer price index for all urban consumers (FRED series \texttt{CPIAUCSL}).} The controls are the six states that form the synthetic Alaska of \citet{JonesMarinescu2022} for the employment rate: Utah, Wyoming, Washington, Nevada, Montana and Minnesota. Oil moves both the economy of Alaska and the income of its government, so the VAR includes the monthly change in the logarithm of the price of West Texas Intermediate crude as a second policy $R_t$, ordered first.\footnote{FRED series \texttt{WTISPLC}, spot price of West Texas Intermediate crude oil.}


The surprise in the dividend must also carry no news about movements of employment specific to Alaska (Assumption \ref{assu:A4}). Until 2015 the formula set the dividend from the income of a fund whose diversified portfolio the economy of Alaska does not move. From 2016 the legislature set it with an eye on a budget that depends on oil revenue, a channel that ordering the price of oil first removes from the innovation.

Table \ref{tab:ak-weights} reports the test of the design and the weights. The outcomes of Alaska behave as integrated series, and the test does not reject a stochastic trend specific to Alaska for any outcome, so I estimate the CVAR in changes (Theorem \ref{thm:diff}). I estimate the weights from the regression of the changes of each Alaska outcome on the changes of the six states, with the dividend and the price of oil as regressors, constrained to be non-negative. The weight spreads over several states for every outcome, and no state receives most of it. The bands account for the estimation of the weights, which I repeat in each bootstrap draw on a moving-block resample of the regression. Table \ref{tab:ak-weights} also reports how much of the main common movement of Alaska the weighted donors reproduce, measured on the first principal component of the changes of the other states (Proposition \ref{prop:wchanges}). The estimated weights fall short of it, as the proposition predicts. Equal weights overshoot it for most outcomes, because Alaska follows the common movement less closely than the average of the six states, so the two weightings miss it in opposite directions.

\safeinput{tab_alaska_weights.tex}

Figure \ref{fig:ak-main} and Table \ref{tab:ak-main} report the responses to a dividend larger by \$1,000 per person. The CVAR finds no fall in the employment rate and no rise in unemployment at any horizon. The VAR without controls finds a large and significant fall in the employment rate, and in the VAR that contains both Alaska and the weighted controls, the controls, which never receive the dividend, fall with it (Proposition \ref{prop:placebo}). The loss of employment that the VAR attributes to the dividend is the national downturn that followed the largest dividends. Participation falls over the two years after the payment, as an income effect on labour supply predicts.\footnote{Supplementary Appendix \ref{sec:app-eb} re-estimates the main specification on the unadjusted series with month effects, with month effects interacted with a linear trend, and with month effects by decade. The absence of a fall in the employment rate and the rise in leisure and hospitality employment hold in every version, while the responses of participation and of total payroll employment depend on how the seasonal pattern of Alaska is removed.}

\begin{figure}[htbp]
\centering\safegraphic[width=\textwidth]{fig_alaska_main.pdf}
\caption{Alaska Permanent Fund Dividend: cumulated response to a dividend larger by \$1,000 per person. CVAR with 68 and 90 percent bands, VAR without controls with its 90 percent band, and response of the weighted controls in the VAR with Alaska and the weighted controls. Rates in percentage points, employment in percent. Monthly, January 1976 (payroll series January 1990) to December 2019.\label{fig:ak-main}}
\end{figure}

Payroll employment rises, and the rise is concentrated in leisure and hospitality, the sector that sells to local customers and where the spending of the dividend lands, while retail trade does not move significantly. For these outcomes the common trends bias the VAR in the opposite direction. After a large dividend the weighted controls lose jobs, and the VAR without controls, which keeps that loss in its response, finds no effect. The common trends therefore produce in the VAR a loss of employment where the dividend has none, and hide the gains in the sector where it has one.

\safeinput{tab_alaska_main.tex}

Because the formula averages the income of five years, a dividend larger than predicted is followed by larger dividends in the next years, and the responses in Table \ref{tab:ak-main} include the effect of those later payments. The impulse causal effect of Theorem \ref{thm:direct} holds the later dividends fixed and isolates the effect of a single payment. Employment in leisure and hospitality peaks within a year of the payment and settles at a lower level, and payroll employment settles after a year.

The estimated weights leave part of the common movement of Alaska in the cycle, and Table \ref{tab:ak-adjusted} in the Supplementary Appendix removes the term that Proposition \ref{prop:wchanges} attributes to it, a multiple of the response of the weighted donors with the same sign. The adjustment brings the responses with estimated weights and with equal weights to the same values, although the two weightings miss the common movement in opposite directions. The term works against the main results. After the adjustment the employment rate still does not fall, and the gains in payroll employment and in leisure and hospitality are larger and remain significant, while the fall in participation loses its significance.


The remaining threat to identification is a component of the Alaska economy that the controls do not share and that moves with the dividend. Table \ref{tab:ak-bounds} in the Supplementary Appendix reports the bounds of Proposition \ref{prop:bounds}. The effects on participation, payroll employment and leisure and hospitality remain significant for biases specific to Alaska that are a substantial fraction of the bias from the common trends, which the response of the controls measures and which, for employment, is itself large. The results also hold when the sample ends in 2015, before the legislature set the amounts, and the point estimates are similar, with wider bands, when it ends in 2007, before the dividend, the rebate and the recession of 2008 (Supplementary Appendix \ref{sec:app-eb}).


\section{Extended unemployment benefits\label{sec:ebsum}}

The second application concerns a policy that reacts to the cycle by design. In the United States, the Extended Benefits programme adds weeks of unemployment insurance when the unemployment of a state crosses thresholds set in federal law, so extensions arrive when unemployment is high and rising. I compare each state in which the programme switched on with the seven states in which it never did, leaving out the controls that border it, using monthly unemployment rates from October 2002 to December 2019. The test of Section \ref{sec:test} rejects no more often than its size would produce, and I estimate the CVAR in changes with equal weights on the controls.

A VAR without controls finds that an extension is followed by a fall in unemployment that keeps growing for three years. The states that never extended benefits show almost the same fall in the VAR that contains both (Proposition \ref{prop:placebo}), so the response of the VAR is the recovery from the recession that triggered the extensions. The CVAR finds a response close to zero at every horizon, with bands several times narrower than those of the VAR without controls, in agreement with the evidence that extensions at long potential durations have little effect on unemployment \citep{AcostaEtAl2025}. Supplementary Appendix \ref{sec:eb} reports the institutions, the data, the discussion of the assumptions, the test of the design, the results state by state, the impulse causal effects, the bounds and the comparison with local-projection difference-in-differences.

\section{The Low Carbon Fuel Standard\label{sec:lcfssum}}

The third application turns to a policy whose cost changes every month. California's Low Carbon Fuel Standard requires refiners and importers of gasoline to cover with credits the carbon intensity of their fuel above a declining benchmark, and whether the cost of the credits reaches consumers is disputed \citep{CARB2024,KleinmanCenter2024}. A comparison of the cost with retail prices mixes it with the movements of crude oil and refining margins that affect every region. I compare the retail price of gasoline in California with the prices of three sets of regions that buy crude oil and refined products in the same markets without a comparable standard, using monthly data from January 2014 to July 2026. The test of Section \ref{sec:test} does not reject a trend specific to California with any set of controls, and I estimate the CVAR in changes with equal weights.

With every set of controls, the CVAR finds that the price of gasoline in California rises by about the increase in the cost within a month and stays near that level, as incidence theory predicts for a cost per gallon borne by the suppliers of a good with inelastic demand. A VAR without controls finds a response that keeps rising well above full pass-through, with much wider bands. Supplementary Appendix \ref{sec:lcfs} reports the institutions, the construction of the cost, the data, the discussion of the assumptions and the robustness of the result across samples, sets of controls and treatments of the pandemic years.

\section{Conclusions\label{sec:conclusions}}

This paper has proposed the control VAR, which uses units not exposed to a policy to remove the common trends of the treated unit and estimates the dynamic effect of the policy on the remaining cycle. The causal framework defines the effect of a policy surprise, including the later reaction of the policy, and the assumptions require controls that share the trends of the treated unit and are not affected by the policy. The theorems show what a recursive VAR identifies, that the response of the cycle removes the bias from the reaction of the policy to the common trends, that the response of the controls measures that bias in any VAR, and that the estimated weights enter inference when the controls are cointegrated among themselves and the policy reacts to the past of the treated unit. The simulations show that the CVAR needs long lags and that a bootstrap that re-estimates the weights delivers the nominal coverage. In the applications, a larger Alaska dividend leaves the employment rate unchanged and raises employment in leisure and hospitality, while a VAR without controls attributes to it a loss of employment that the states without the dividend also show. Extended unemployment benefits have a small average effect on unemployment, while a VAR without controls attributes to them large reductions in the states where they coincided with the deepest recessions. The cost of the Low Carbon Fuel Standard passes almost fully to the price of gasoline within a month. A natural extension is a panel version of the CVAR, in which several treated units share the same controls.

\newpage
\appendix
\section{Proofs of the main results\label{sec:proofs-main}}

\begin{lem}[Responses of distributed lags\label{lem:dl}]
Let $u_t$ be orthogonal to every element of $\mathcal F_{t-1}$, with $\E[u_t^2]>0$, let $W_t$ be covariance stationary and measurable with respect to $\mathcal F_t$, let $\psi_k=\E[W_{t+k}u_t]/\E[u_t^2]$, and let $\sum_{j\geq0}|\tau_j|<\infty$. Then, for every $h\geq0$,
\[
\frac{\E\big[\sum_{j\geq0}\tau_jW_{t+h-j}\,u_t\big]}{\E[u_t^2]}=\sum_{j=0}^{h}\tau_j\psi_{h-j},
\qquad
\E\Big[\sum_{j\geq0}\tau_jW_{t-1-j}\,u_t\Big]=0 .
\]
\end{lem}

\paragraph{Proof of Lemma \ref{lem:dl}.} Since $\sum_j|\tau_j|\,\E|W_{t+h-j}u_t|\leq\sum_j|\tau_j|(\E[W_t^2]\E[u_t^2])^{1/2}<\infty$, the sum and the expectation can be interchanged. For $j>h$, and for every term of the second sum, the policy is dated $t-1$ or earlier, so it belongs to $\mathcal F_{t-1}$ and is orthogonal to $u_t$. For $j\leq h$, $\E[W_{t+h-j}u_t]=\psi_{h-j}\E[u_t^2]$. $\square$

\paragraph{Proof of Theorem \ref{thm:var}.} By Assumptions \ref{assu:A1} to \ref{assu:A3}, $Y^1_{t+h}-Y^1_{t-1}=\lambda'(F_{t+h}-F_{t-1})+(e_{t+h}-e_{t-1})+\sum_{j\geq0}\tau_j(W_{t+h-j}-W_{t-1-j})$. Multiply by $\tilde u_t$, take expectations and divide by $\E[\tilde u_t^2]$. By Lemma \ref{lem:dl}, the policy terms give $\sum_{j=0}^h\tau_j\tilde\psi_{h-j}$, which is $\tilde\theta_h$ by \eqref{eq:decomp}. The expectations are finite by Assumption \ref{assu:A5}. $\square$

\paragraph{Proof of Theorem \ref{thm:cycle}.} By Assumption \ref{assu:A2} and $\lambda=\Lambda'b$, the common trends cancel in $Y^1_t(0)-b'Y^0_t=\lambda'F_t+e_t-b'\Lambda F_t-b'v_t=e_t-b'v_t=G^0_t$; by Assumption \ref{assu:A1} the controls do not depend on the policy; and by Assumption \ref{assu:A3}, $G_t=G^0_t+\sum_{j\geq0}\tau_jW_{t-j}$. Multiplying $G_{t+h}$ by $u_t$ and applying Lemma \ref{lem:dl} gives \eqref{eq:cyclebias}, with the policy terms equal to $\theta_h$ by \eqref{eq:decomp}.

For the second statement, let $Z_t$ follow a VAR, possibly of infinite order. The response of $G_{t+h}$ to the recursive innovation of $W_t$, normalised to a unit impact, is the coefficient on $u_t$ in the projection of $G_{t+h}$ on $u_t$, $R_t$ and $\mathcal F_{t-1}$. Since $u_t$ is orthogonal to $R_t$ and to $\mathcal F_{t-1}$, the coefficient is $\E[G_{t+h}u_t]/\E[u_t^2]$ \citep[Proposition 1]{PlagborgMollerWolf2021}. $\square$

\paragraph{Proof of Theorem \ref{thm:ident}.} By Theorem \ref{thm:cycle}, $\gamma_h-\theta_h=\E[G^0_{t+h}u_t]/\E[u_t^2]$, which is zero for every $h\geq0$ by Assumption \ref{assu:A4}. The second equality is the decomposition \eqref{eq:decomp}, which holds under Assumption \ref{assu:A3}. $\square$

\paragraph{Proof of Theorem \ref{thm:direct}.} (a) By Theorem \ref{thm:ident} and \eqref{eq:decomp}, $\gamma_h=\sum_{j=0}^h\tau_j\psi_{h-j}$ for every $h\geq0$, with $\psi_0=1$ because $u_t$ is the part of $W_t$ orthogonal to $R_t$ and $\mathcal F_{t-1}$. The system is triangular with unit diagonal, so $\tau_0=\gamma_0$ and $\tau_h=\gamma_h-\sum_{j=0}^{h-1}\tau_j\psi_{h-j}$ for $h\geq1$. Both $\gamma_h$ and $\psi_h$ are responses of the VAR in $Z_t$ to the innovation of the policy, by the second part of the proof of Theorem \ref{thm:cycle}, applied to $G_t$ and to $W_t$.

(b) Let $\kappa_k=\E[W_{t+k}(0)u_t]/\E[u_t^2]$ be the response of the path the policy would have followed without the surprise, with $\kappa_0=0$. Under Assumption \ref{assu:A3b}, $W_{t+k}=W_{t+k}(0)+\phi_kW_t$, so $\psi_k=\phi_k+\kappa_k$. Substituting in \eqref{eq:decomp} gives \eqref{eq:pathbias}. Under Assumptions \ref{assu:A3} and \ref{assu:A3b}, $Y_{t+h}\langle w\rangle-Y_{t+h}\langle 0\rangle=\sum_{j=0}^h\tau_j\big(W_{t+h-j}(w)-W_{t+h-j}(0)\big)=w\sum_{j=0}^h\tau_j\phi_{h-j}=w\,\theta^\ast_h$. If $\kappa_k=0$ for all $k\geq1$, then $\theta_h=\theta^\ast_h$, and $\gamma_h=\theta^\ast_h$ by Theorem \ref{thm:ident}. $\square$

\paragraph{Proof of Proposition \ref{prop:placebo}.} (a) Since $Y^1_t=G_t+C_t$, the identity holds for the moments, after division by $\E[\breve u_t^2]$. It holds for the Cholesky responses as well, because the VARs in $(R_t',W_t,Y^1_t,C_t)'$ and in $(R_t',W_t,G_t,C_t)'$ are linear transformations of each other that leave the first two variables unchanged, so they have the same policy innovation. Since $G_{t-1}$ belongs to the information set of $\breve u_t$, $\gamma^G_h=\E[G_{t+h}\breve u_t]/\E[\breve u_t^2]$.

(b) By Assumption \ref{assu:A2}, $C_t=b'\Lambda F_t+b'v_t=\lambda'F_t+b'v_t$, and by Assumption \ref{assu:A1} it does not depend on the policy. Substituting in $\gamma^C_h$ gives the expression, which is finite by Assumption \ref{assu:A5}; its second term is zero if $\E[v_{t+j}\breve u_t]=0$ for all $j$.

(c) The proofs of Theorems \ref{thm:cycle} and \ref{thm:ident} apply with $\breve u_t$ in place of $u_t$, because the past of $W_t$ belongs to the information set of $\breve u_t$. $\square$

\paragraph{Proof of Proposition \ref{prop:bounds}.} By Theorem \ref{thm:cycle} applied with the innovation $\breve u_t$, $\gamma^G_h=\breve\theta_h+\breve B_h$, so $\breve\theta_h=\gamma^G_h-\breve B_h$, and Assumption \ref{assu:A4b} gives (a). By Proposition \ref{prop:placebo}(a), $\gamma^{Y^1}_h-\gamma^G_h=\gamma^C_h$, so $|\gamma^{Y^1}_h-\gamma^G_h|=|\gamma^C_h|$, which is at most $\lambda|\gamma^C_h|$ only if $\lambda\geq1$ or $\gamma^C_h=0$. $\square$

\paragraph{Proof of Theorem \ref{thm:est}.} Write $\eta_1=\hat c_1-c_1^\ast$ and $\eta_2=\hat c_2-c_2^\ast$ for the errors in the weights on the trends and on the stationary combinations, and let a tilde denote a series demeaned over the sample, so that $\tilde z_t$ is the demeaned vector $(R_t',W_t,G^\ast_t)'$. Let $M$ collect the population autocovariances of $\tilde z_t$ up to lag $p$, $\tilde M$ their sample counterparts computed with the cycle $G^\ast_t$, and $\hat M$ the same sample autocovariances computed with the estimated cycle $\hat G_t$. Let $\mathcal Q$ be the probability limit of $T^{-1}\sum_t\tilde S_t\tilde S_t'$ after the projection on the other stationary regressors of \eqref{eq:dols}, and $\varepsilon^\infty_t$ the error of \eqref{eq:dols} with all leads and lags.

(a) In the rotated coordinates, \eqref{eq:dols} is a cointegrating regression of $Y^1_t$ on the $r$ integrated regressors $X_{1t}$, whose trends have full rank, and on stationary regressors: $S_t$, the policy terms, the leads and lags of $\Delta X_{1t}$ and the constant. The leads and lags make the error asymptotically orthogonal to the innovations of $X_{1t}$, so the results for dynamic least squares \citep{Saikkonen1991,StockWatson1993}, applied to the canonical form of \citet{SimsStockWatson1990}, give $T\eta_1=O_p(1)$, and $\sqrt{T}$-consistency and asymptotic normality of the coefficients on the stationary regressors, which converge to the population projection coefficients $(a^\ast,c_2^\ast)$; in particular $\sqrt{T}\eta_2=\mathcal Q^{-1}T^{-1/2}\sum_t\tilde S_t\varepsilon^\infty_t+o_p(1)$ without the constraint. Under the constraint $b\geq0$, which is a convex cone in the coordinates $c$ as well, the scaled estimation error converges to the projection of the unconstrained limit on the tangent cone at $c^\ast$ \citep{Andrews1999}. This limit is $O_p(1)$ in both blocks, and it is normal when $c^\ast$ is in the interior.\footnote{With $L$ fixed, least squares remains super-consistent for the coefficients of the integrated regressors \citep{Stock1987}, so $T\eta_1=O_p(1)$ still holds, and parts (a) and (b) hold with $b^\ast$ replaced by the minimiser of the population criterion with $L$ leads and lags.}

(b) The argument has four steps.

Step 1 (smoothness). The least-squares estimator of the VAR and its residual covariance are smooth functions of the sample autocovariances of $\tilde z_t$ up to lag $p$, so $\hat\gamma_h=g_h(\hat M)$ and $\tilde\gamma_h=g_h(\tilde M)$ for a function $g_h$ that is continuously differentiable at $M$.

Step 2 (the effect of the estimated weights on the moments). Since $\hat G_t-G^\ast_t=-(\hat a-a^\ast)-\eta_1'X_{1t}-\eta_2'S_t$, and the constant disappears in the demeaned series,
\[
\tilde{\hat G}_t=\tilde G^\ast_t-\eta_1'\tilde X_{1t}-\eta_2'\tilde S_t .
\]
Two kinds of entries of $\hat M$ change. The first are cross moments of the cycle with another variable of the system, an element $\varsigma_t$ of $\tilde z_{t-k}$ other than the cycle:
\[
\frac1T\sum_t\tilde{\hat G}_t\varsigma_t-\frac1T\sum_t\tilde G^\ast_t\varsigma_t=-\eta_1'\,\frac1T\sum_t\tilde X_{1t}\varsigma_t-\eta_2'\,\frac1T\sum_t\tilde S_t\varsigma_t .
\]
The sample covariance between a demeaned integrated process and a stationary one is $O_p(1)$ \citep{Phillips1988}, so the first term is $O_p(T^{-1})\,O_p(1)=O_p(T^{-1})$. By the law of large numbers, $T^{-1}\sum_t\tilde S_t\varsigma_t=\E[\tilde S_t\varsigma_t]+O_p(T^{-1/2})$, so the second term is $-\eta_2'\E[\tilde S_t\varsigma_t]+O_p(T^{-1})$.

The second kind are the moments of the cycle with itself, at lags $k=0,\dots,p$. Expanding the product,
\begin{align*}
\frac1T\sum_t\tilde{\hat G}_t\tilde{\hat G}_{t-k}-\frac1T\sum_t\tilde G^\ast_t\tilde G^\ast_{t-k}
={}&-\eta_1'\,\frac1T\sum_t\big(\tilde X_{1t}\tilde G^\ast_{t-k}+\tilde X_{1,t-k}\tilde G^\ast_t\big)
-\eta_2'\,\frac1T\sum_t\big(\tilde S_t\tilde G^\ast_{t-k}+\tilde S_{t-k}\tilde G^\ast_t\big)\\
&+\eta_1'\Big(\frac1T\sum_t\tilde X_{1t}\tilde X_{1,t-k}'\Big)\eta_1
+\eta_1'\Big(\frac1T\sum_t\big(\tilde X_{1t}\tilde S_{t-k}'+\tilde X_{1,t-k}\tilde S_t'\big)\Big)\eta_2\\
&+\eta_2'\Big(\frac1T\sum_t\tilde S_t\tilde S_{t-k}'\Big)\eta_2 .
\end{align*}
The orders of the five terms, with $\eta_1=O_p(T^{-1})$ and $\eta_2=O_p(T^{-1/2})$ from part (a), are as follows. The first term is $O_p(T^{-1})\,O_p(1)$, by the same covariance result as above. The second term equals $-\eta_2'\big(\E[\tilde S_t\tilde G^\ast_{t-k}]+\E[\tilde S_{t-k}\tilde G^\ast_t]\big)+O_p(T^{-1})$, by the law of large numbers. In the third term, the sample second moment of a demeaned integrated process grows with the sample, $T^{-1}\sum_t\tilde X_{1t}\tilde X_{1,t-k}'=O_p(T)$, because $T^{-2}\sum_t\tilde X_{1t}\tilde X_{1,t-k}'$ converges to the integral of the outer product of a demeaned Brownian motion \citep{Phillips1987}. The third term is therefore $O_p(T^{-1})\,O_p(T)\,O_p(T^{-1})=O_p(T^{-1})$: the super-consistency of the weights on the trends exactly offsets the growth of the trends. The fourth term is $O_p(T^{-1})\,O_p(1)\,O_p(T^{-1/2})=O_p(T^{-3/2})$, and the fifth is $O_p(T^{-1/2})\,O_p(1)\,O_p(T^{-1/2})=O_p(T^{-1})$. Only the second term is of order $T^{-1/2}$.

Collecting both kinds of entries,
\[
\sqrt{T}(\hat M-\tilde M)=\mathcal C\sqrt{T}\,\eta_2+O_p(T^{-1/2}),
\]
where $\mathcal C$ is a fixed matrix whose rows are minus the population covariances of $\tilde S_t$ with the elements of $\tilde z_t$ at the leads and lags that enter $M$.

Step 3 (delta method). By Step 1 and Step 2,
\[
\sqrt{T}(\hat\gamma_h-\tilde\gamma_h)=\nabla g_h(M)'\mathcal C\sqrt{T}\,\eta_2+o_p(1),
\]
which is \eqref{eq:firststep} with $J_h=\mathcal C'\nabla g_h(M)$. Equivalently, $J_h$ is the derivative at $\delta=0$ of the population response of the cycle $G^\ast_t-\delta'S_t$ to the policy innovation.

Step 4 (joint normality). Up to $o_p(1)$, $\sqrt{T}(\tilde\gamma_h-\gamma_h)$ is a linear function of $\sqrt{T}(\tilde M-M)$, and $\sqrt{T}\,\eta_2=\mathcal Q^{-1}T^{-1/2}\sum_t\tilde S_t\varepsilon^\infty_t+o_p(1)$ by part (a). Both are normalised sums of products of the elements of $(R_t',W_t,G^\ast_t,\Delta X_{1t}',S_t')'$ and of the error of \eqref{eq:dols} at lags up to $p$. These products are functions of at most $p+1$ consecutive observations, so they are strong mixing with coefficients $\alpha(m-p)$ for $m>p$, and they have finite moments of order $2+\epsilon$. The central limit theorem for strong mixing sequences \citep{Herrndorf1984}, with the Cram\'er-Wold device, gives joint asymptotic normality. Proposition \ref{prop:J} gives $J_h$ and conditions under which it is zero.

(c) If $r=K$, there are no stationary combinations of the controls, $S_t$ is absent, and only the terms in $\eta_1$ remain in Step 2. They are $O_p(T^{-1})=o_p(T^{-1/2})$, and the argument uses only $\eta_1=O_p(T^{-1})$, so it holds with and without the constraint. The case $p=p_T\to\infty$ is part (a) of Theorem \ref{thm:joint}. $\square$

\paragraph{Proof of Theorem \ref{thm:joint}.} I keep the notation of the proof of Theorem \ref{thm:est}, with $\tilde z_t$ the demeaned vector of the VAR and $\mathcal Q$ the limit of the second moments of the stationary combinations of the controls, and write $\hat{\mathcal Q}$ for its bootstrap counterpart. For a VAR($p$), $\tilde A(p)$ denotes the least-squares coefficients computed with the cycle $G^\ast_t$, $\mathcal C_p$ the matrix of the population covariances of $\tilde S_{t-l}$ with the elements of $\tilde z_{t-k}$, $l,k\leq p$, and $J_h(p)$ the derivative of the population response of the VAR($p$) with respect to the weights on $S_t$. I use four results on autoregressions of growing order. For a stationary process with a VAR($\infty$) representation satisfying the conditions of the theorem, let $\hat\Gamma_p$ be the sample second-moment matrix of the regressors of a VAR($p$) and $\Gamma_p$ its population counterpart. (LR1) $\|\hat\Gamma_p^{-1}-\Gamma_p^{-1}\|=O_p(pT^{-1/2})$ in operator norm, and $\|\Gamma_p^{-1}\|$ is bounded uniformly in $p$. (LR2) The least-squares coefficients satisfy $\|\hat A(p)-A(p)\|=O_p(p^{1/2}T^{-1/2})$, where $A(p)$ collects the first $p$ coefficients of the VAR($\infty$). (LR3) For any sequence of non-random vectors $\ell_p$ with bounded norm, $\sqrt{T}\ell_p'\mathrm{vec}(\hat A(p)-A(p))=T^{-1/2}\sum_t\ell_p'\mathrm{vec}(x_t(p)\epsilon_t')\Gamma_p^{-1}+o_p(1)$, which is asymptotically normal \citep{LewisReinsel1985}. (IK) The residual bootstrap of the VAR($p$) reproduces (LR1) to (LR3) in probability, and for smooth functions of the slope coefficients and of the innovation covariance it reproduces the limit distribution of the estimator \citep{InoueKilian2002}.

(a) The estimator of the weights does not involve $p$, so Theorem \ref{thm:est}(a) holds unchanged. Write $\hat G_t-G^\ast_t=-(\hat a-a^\ast)-(\hat c_1-c_1^\ast)'X_{1t}-(\hat c_2-c_2^\ast)'S_t$; the constant is absorbed by the constant of the VAR. Replacing $G^\ast_t$ by $\hat G_t$ changes the sample second moments of the regressors and their cross moments with the dependent variables. The term in $X_{1t}$ changes each entry by $(\hat c_1-c_1^\ast)'T^{-1}\sum_t\tilde X_{1,t-l}\tilde z_{t-k}=O_p(T^{-1})\,O_p(1)$, because the sample covariance between a demeaned integrated process and a stationary one is $O_p(1)$ \citep{Phillips1988}, and its quadratic term is $O_p(T^{-1})$. Over the $p$ lags the change in $\hat\Gamma_p$ and in the cross moments is $O_p(pT^{-1})$ in Frobenius norm. The term in $S_t$ changes each entry by $(\hat c_2-c_2^\ast)'\big(\E[\tilde S_{t-l}\tilde z_{t-k}]+O_p(T^{-1/2})\big)$, so the change is $(\hat c_2-c_2^\ast)'\mathcal C_p+O_p(p^{1/2}T^{-1})$, and $\mathcal C_p$ has bounded norm. By (LR1), the coefficients of the VAR in $\hat Z_t$ then satisfy
\[
\hat A(p)-\tilde A(p)=\Gamma_p^{-1}\big[(\hat c_2-c_2^\ast)'\mathcal C_p\big]^{\dagger}+O_p(pT^{-1})+O_p(p^{3/2}T^{-1}),
\]
where $[\cdot]^\dagger$ is the linear map from a perturbation of the moments to the perturbation of the coefficients, evaluated at the population moments. The remainders are $o_p(T^{-1/2})$ because $p^3/T\to0$. The response at a fixed horizon $h$ depends on the first $h$ coefficient matrices and on the innovation covariance through a function that is continuously differentiable with bounded derivatives, so
\[
\sqrt{T}(\hat\gamma_h-\gamma_h)=\sqrt{T}(\tilde\gamma_h-\gamma_h)+J_h(p)'\sqrt{T}(\hat c_2-c_2^\ast)+o_p(1).
\] As $p\to\infty$, the population projections of the VAR($p$) converge to those of the VAR($\infty$), because $\sum_l l\|\Phi_l\|<\infty$, and $J_h(p)\to J_h$, which gives \eqref{eq:firststep}. By (LR3) and the delta method, $\sqrt{T}(\tilde\gamma_h-\gamma_h)$ equals a normalised sum $T^{-1/2}\sum_t\xi_{t,T}$ with $\xi_{t,T}$ a linear function of $\mathrm{vec}(x_t(p)\epsilon_t')$ and of $\epsilon_t\epsilon_t'-\Sigma$, plus $o_p(1)$. By the proof of Theorem \ref{thm:est}(a), $\sqrt{T}(\hat c_2-c_2^\ast)=\mathcal Q^{-1}T^{-1/2}\sum_t\tilde S_t\varepsilon^\infty_t+o_p(1)$. The two sums are sums of a strong mixing triangular array with uniformly bounded moments of order $2+\epsilon$, so the Cram\'er-Wold device and the central limit theorem for strong mixing arrays give their joint asymptotic normality.

(b) Step 1 (the bootstrap world). In the bootstrap world, $Q^\ast_t$ follows the VAR($p$) estimated in step (i), with innovations drawn from the centred residuals. Its coefficients and innovation covariance converge in probability to those of the VAR($p$) approximation of $Q_t$, and hence to those of its VAR($\infty$) representation, by (LR2) and the summability condition. Let $b^\natural$ be the weights to which \eqref{eq:dols} converges in the bootstrap world, and $\gamma^\ast_h$ the response to which the CVAR converges there. Both are the population pseudo-true values of part (a) computed for the process $Q^\ast_t$, so they are continuous functions of its coefficients and innovation covariance, and they converge in probability to $b^\ast$ and $\gamma_h$. Step (v) computes $\gamma^\ast_h$ up to a simulation error that vanishes as the length of the long path grows.

Step 2 (bootstrap expansion). Since $Y^{1\ast}_t-\hat a-\hat b'Y^{0\ast}_t=G^\ast_t$, the bootstrap estimation error of the weights relative to $b^\natural$ is the error of the regression \eqref{eq:dols} of $G^\ast_t$ on the bootstrap regressors, centred at its bootstrap pseudo-true value. The controls $Y^{0\ast}_t$ have the rotated components $X^\ast_{1t}$, integrated, and $S^\ast_t$, stationary, generated jointly with $(R^\ast_t,W^\ast_t,G^\ast_t)$, so the proof of Theorem \ref{thm:est}(a), applied conditionally on the data with (IK) in place of (LR1) to (LR3), gives $\hat c^\ast_1-c^\natural_1=O_{p^\ast}(T^{-1})$ and $\sqrt{T}(\hat c^\ast_2-c^\natural_2)=\hat{\mathcal Q}^{-1}T^{-1/2}\sum_t\tilde S^\ast_t\varepsilon^\ast_t+o_{p^\ast}(1)$, in probability. Part (a), applied in the same way, gives
\[
\sqrt{T}(\hat\gamma^\ast_h-\gamma^\ast_h)=\sqrt{T}(\tilde\gamma^\ast_h-\gamma^\ast_h)+J^\ast_h(p)'\sqrt{T}(\hat c^\ast_2-c^\natural_2)+o_{p^\ast}(1),
\]
where $\tilde\gamma^\ast_h$ is the bootstrap response with the weights $b^\natural$ and $J^\ast_h(p)$ is the derivative of the response with respect to the weights on $S^\ast_t$ in the bootstrap world. By Step 1, $J^\ast_h(p)\to J_h$ in probability.

Step 3 (joint limit). In both worlds the pair formed by the response with the pseudo-true weights and the scaled error of the weights on $S_t$ is, up to negligible terms, a fixed linear functional of the scaled sample autocovariances of $Q_t$ (respectively $Q^\ast_t$) at lags up to $p$, centred at their population values in that world. For a linear process with i.i.d. innovations, the limit covariance of such a functional is a continuous function of the coefficients of the process, of the innovation covariance and of the fourth cumulants of the innovations. In the bootstrap world these are the coefficients of the fitted VAR($p$), the residual covariance and the fourth cumulants of the centred residuals, which converge in probability to their population values \citep{InoueKilian2002}. The bootstrap limit is therefore the normal limit of part (a).

Step 4 (coverage). The shifted draws of step (v) are $\hat\gamma_h+(\hat\gamma^\ast_h-\gamma^\ast_h)$, so their quantiles are $\hat\gamma_h$ plus the quantiles of $\hat\gamma^\ast_h-\gamma^\ast_h$. The limit of $\sqrt{T}(\hat\gamma_h-\gamma_h)$ is normal with mean zero, hence symmetric, and continuous, so the convergence of the distribution functions is uniform and the percentile band of the shifted draws has asymptotic coverage $1-\alpha$. The estimated rank and $\hat D_1$ are held fixed in both worlds; the rank is correct with probability approaching one, and $\hat D_1$ converges at the rate $T^{-1}$, so neither affects the limits. $\square$

\paragraph{Proofs for Section \ref{sec:test}.} The proof of Theorem \ref{thm:po} has three steps. The weights and the policy coefficients of the quasi-differenced regression are $\sqrt{T}$-consistent under the null and under local alternatives (Lemma \ref{lem:orders}). Their estimation changes $S(\bar\alpha)$ and $S(1)$ by amounts of order one, which cancel in the numerator of $P_T$ up to $o_p(1)$, because $\bar\alpha-1=\bar c/T$ (proof of Theorem \ref{thm:po}). It leaves the estimator of the long-run variance consistent (Lemma \ref{lem:lrv}). Lemma \ref{lem:fclt} collects the limits that the three steps use; their bootstrap counterparts, and the proof of Proposition \ref{prop:sieve}, are in Supplementary Appendix \ref{sec:supp-testboot}. Throughout, $\Rightarrow$ denotes weak convergence, $\|\cdot\|$ the Euclidean norm (Frobenius for matrices), and all limits are as $T\to\infty$ with $c$ and $\bar c$ fixed. The statistic is invariant to nonsingular linear transformations of the regressors and of the controls, so I rotate the controls into the $r$ trend directions $x_t=D_1'Y^0_t$ and the $K-r$ stationary combinations $s_t=H'Y^0_t$, and collect the stationary regressors in $w_t=(s_t',P_t')'$ and the stochastic regressors in $V_t=(x_t',w_t')'$. For any variable $z_t$, $z^\alpha_1=z_1$ and $z^\alpha_t=z_t-\alpha z_{t-1}$ for $t\geq2$; $d^\alpha$ is the quasi-differenced constant and $M_\alpha=I-d^\alpha(d^{\alpha\prime}d^\alpha)^{-1}d^{\alpha\prime}$. Under \eqref{eq:testmodel}, $Y^{1\alpha}=\mu d^\alpha+V^\alpha\vartheta+u^\alpha$, with $\vartheta$ collecting $b$ and $\delta(L)$ in the rotated coordinates, so that
\begin{equation}
S(\alpha)=u^{\alpha\prime}M_\alpha u^\alpha-A_\alpha'Q_\alpha^{-1}A_\alpha,\qquad Q_\alpha=V^{\alpha\prime}M_\alpha V^\alpha,\quad A_\alpha=V^{\alpha\prime}M_\alpha u^\alpha,
\label{eq:Sdecomp}
\end{equation}
where the first term, $S_u(\alpha)$, is the sum of squares of the GLS regression of $u_t$ on the constant alone. The test regression is unconstrained, so no boundary issue arises: the restriction $b\geq0$ belongs to the estimator of Section \ref{sec:theorems} and plays no role in the test.

I write the linear processes of Assumption \ref{assu:T} as $v_t=\psi(L)\varepsilon_t$ and $\Delta V_t-E\Delta V_t=\Phi(L)\eta_t$, where $\{(\varepsilon_t,\eta_t')'\}$ is independent and identically distributed with mean zero, finite fourth moments and covariance $\mathrm{diag}(\sigma^2_\varepsilon,\Sigma_\eta)$, the off-diagonal block being zero because $v_t$ and $\Delta V_t$ are uncorrelated at all leads and lags, and $\sum_jj(|\psi_j|+\|\Phi_j\|)<\infty$. Then $\omega=\psi(1)\sigma_\varepsilon>0$ and $\Gamma=E[\Delta V_t\Delta V_t']$ is positive definite.

\begin{lem}\label{lem:fclt}
Under Assumption \ref{assu:T} and $\rho_T=1+c/T$:
\begin{enumerate}
\item[(a)] on $[0,1]$, \[\big(T^{-1/2}u_{[sT]},\,T^{-1/2}x_{[sT]}'\big)\Rightarrow\big(\omega W_c(s),\,B_x(s)'\big),\] where $W_c(s)=\int_0^se^{c(s-r)}dW_0(r)$, $W_0$ is a standard Brownian motion and $B_x$ is a Brownian motion independent of $W_0$;
\item[(b)] $T^{-1}\sum_t\Delta V_t\Delta V_t'\to_p\Gamma$ and $T^{-2}\sum_tx_{t-1}x_{t-1}'\Rightarrow\int_0^1B_xB_x'$;
\item[(c)] $\sum_t\Delta V_tv_t=O_p(T^{1/2})$; $T^{-1}\sum_t\Delta V_tu_{t-1}\Rightarrow\omega\int_0^1W_c\,dB_V$, where $B_V$ is the Brownian motion of the partial sums of $\Delta V_t$; $T^{-1}\sum_tx_{t-1}v_t\Rightarrow\omega\int_0^1B_x\,dW_0$; and $T^{-1}\sum_tw_{t-1}v_t=O_p(T^{-1/2})$.
\end{enumerate}
\end{lem}

\noindent Proof. (a) By the Beveridge-Nelson decomposition $v_t=\psi(1)\varepsilon_t+\tilde v_{t-1}-\tilde v_t$ with $\tilde v_t$ stationary, so $\max_{t\leq T}|T^{-1/2}\sum_{s\leq t}v_s-\psi(1)T^{-1/2}\sum_{s\leq t}\varepsilon_s|=o_p(1)$, and the same holds for $\Delta x_t$ with $\Phi_x(1)$ \citep[Theorem 3.4]{PhillipsSolo1992}. The partial sums of $(\varepsilon_t,\eta_t')'$ satisfy a multivariate invariance principle with a block-diagonal covariance, so the limits are jointly Gaussian with zero cross-covariance, hence independent. Writing $u_t=\sum_{s=1}^t\rho_T^{t-s}v_s+\rho_T^tu_0$, the near-integrated limit $T^{-1/2}u_{[sT]}\Rightarrow\omega W_c(s)$ follows from \citet[Lemma 1]{Phillips1987}, and $\rho_T^tT^{-1/2}u_0\to_p0$ uniformly in $t$ by Assumption \ref{assu:T}(iii).

(b) The first limit is the weak law of large numbers for sample second moments of linear processes with absolutely summable coefficients \citep[Theorem 3.7, Remark 3.9(i) and Section 4(i)]{PhillipsSolo1992}; the second follows from (a) and the continuous mapping theorem.

(c) For the first statement, each element of $\sum_t(\Delta V_t-E\Delta V_t)v_t$ is a sample cross-covariance at lag zero of two elements of the vector linear process $(v_t,\Delta V_t')'$, whose innovations are i.i.d. with finite fourth moments and whose coefficients satisfy $\sum_jj(|\psi_j|+\|\Phi_j\|)<\infty$. Its population value is zero, because $v_t$ and $\Delta V_t$ are uncorrelated at all leads and lags, and the central limit theorem for sample covariances of linear processes \citep[Theorem 3.8, Remark 3.9(i) and Section 4(i)]{PhillipsSolo1992} makes the normalised sum $O_p(1)$, so the sum is $O_p(T^{1/2})$; the mean of $\Delta V_t$ multiplies $\sum_tv_t=O_p(T^{1/2})$. The second and third statements are the convergence of sample covariances between a near-integrated process and a stationary linear process to stochastic integrals \citep{Phillips1987,Phillips1988}; the one-sided long-run covariance corrections $\sum_{j\geq1}E[v_{t-j}\Delta V_t']$ and $\sum_{j\geq0}E[\Delta x_{t-j}v_t]$ vanish because the cross-covariances are zero at all lags. The last statement follows as the first, since $w_{t-1}$ is stationary. $\square$

\begin{lem}\label{lem:orders}
Under the conditions of Lemma \ref{lem:fclt}, $T^{-1}Q_1\to_p\Gamma$, $Q_{\bar\alpha}-Q_1=O_p(1)$, $A_1=O_p(T^{1/2})$ and $A_{\bar\alpha}-A_1=O_p(1)$. Consequently $\hat\vartheta(\bar\alpha)-\vartheta=Q_{\bar\alpha}^{-1}A_{\bar\alpha}=O_p(T^{-1/2})$.
\end{lem}

\noindent Proof. For $t\geq2$, $V^{\bar\alpha}_t=\Delta V_t-(\bar c/T)V_{t-1}$ and $u^{\bar\alpha}_t=\Delta u_t-(\bar c/T)u_{t-1}$, with $\Delta u_t=v_t+(c/T)u_{t-1}$. Before the projection on the constant,
\[
\textstyle\sum_{t\geq2}V^{\bar\alpha}_tV^{\bar\alpha\prime}_t-\sum_{t\geq2}\Delta V_t\Delta V_t'=-\frac{\bar c}{T}\sum(\Delta V_tV_{t-1}'+V_{t-1}\Delta V_t')+\frac{\bar c^2}{T^2}\sum V_{t-1}V_{t-1}'.
\]
By Lemma \ref{lem:fclt}, the blocks of the right-hand side involving $x_{t-1}$ are $O_p(1)$: $T^{-1}\sum\Delta V_tx_{t-1}'$ converges to a stochastic integral plus a long-run covariance and $T^{-2}\sum x_{t-1}x_{t-1}'$ to $\int B_xB_x'$; the blocks involving only $w_{t-1}$ are $O_p(T^{-1/2})$ and $O_p(T^{-1})$. Similarly
\[
\textstyle\sum_{t\geq2}V^{\bar\alpha}_tu^{\bar\alpha}_t-\sum_{t\geq2}\Delta V_t\Delta u_t=-\frac{\bar c}{T}\sum(\Delta V_tu_{t-1}+V_{t-1}\Delta u_t)+\frac{\bar c^2}{T^2}\sum V_{t-1}u_{t-1}=O_p(1),
\]
using $T^{-1}\sum V_{t-1}\Delta u_t=T^{-1}\sum V_{t-1}v_t+cT^{-2}\sum V_{t-1}u_{t-1}=O_p(1)$ and $T^{-2}\sum x_{t-1}u_{t-1}\Rightarrow\omega\int B_xW_c$. Finally $\sum\Delta V_t\Delta u_t=\sum\Delta V_tv_t+(c/T)\sum\Delta V_tu_{t-1}=O_p(T^{1/2})$ and $T^{-1}\sum_{t\geq2}\Delta V_t\Delta V_t'\to_p\Gamma$. The first observation contributes $V_1V_1'$ and $V_1u_1$, which are $O_p(1)$ for both values of $\alpha$. The projection on the quasi-differenced constant subtracts $(d^{\alpha\prime}d^\alpha)^{-1}(d^{\alpha\prime}V^\alpha)'(d^{\alpha\prime}u^\alpha)$ and the analogous term from $Q_\alpha$. For $\alpha=1$ it removes only the first observation. For $\alpha=\bar\alpha$, $d^{\bar\alpha\prime}d^{\bar\alpha}=1+(T-1)\bar c^2/T^2\to1$, $d^{\bar\alpha\prime}V^{\bar\alpha}=V_1-(\bar c/T)\sum_{t\geq2}V^{\bar\alpha}_t=V_1-(\bar c/T)(V_T-V_1)+O_p(1)=O_p(1)$, and in the same way $d^{\bar\alpha\prime}u^{\bar\alpha}=O_p(1)$. Both corrections are therefore $O_p(1)$, which proves the four orders; the last statement follows from $Q_{\bar\alpha}^{-1}=O_p(T^{-1})$ and $A_{\bar\alpha}=O_p(T^{1/2})$. $\square$

\begin{lem}\label{lem:lrv}
Let $\hat\omega^2(k)=\hat s^2_k/(1-\sum_{j=1}^k\hat\beta_{j,k})^2$, where $\hat s^2_k$ and $\hat\beta_{j,k}$ are the residual variance and lag coefficients of the regression of $\Delta\hat u_t$ on $\hat u_{t-1}$ and $\Delta\hat u_{t-1},\dots,\Delta\hat u_{t-k}$, and $\hat u_t$ is the residual in levels of the regression with $\alpha=\bar\alpha$; let $\tilde\omega^2(k)$ be the same estimator computed from the GLS-demeaned cycle $\tilde u_t$. Under the conditions of Lemma \ref{lem:fclt} and $k_T\to\infty$ with $k_T^3/T\to0$, $\max_{k\leq k_T}|\hat\omega^2(k)-\tilde\omega^2(k)|=o_p(1)$. If, in addition, the lag order satisfies $\mathbb P(\underline k_T\leq\hat k\leq k_T)\to1$, where $\underline k_T\to\infty$ with $T^{1/2}\sum_{j>\underline k_T}|\beta_j|\to0$ for the coefficients $\beta_j$ of the autoregressive representation of $v_t$, or $\underline k_T$ is at least the order of that representation when it is finite, then $\hat\omega^2(\hat k)\to_p\omega^2$ under the null and under local alternatives.
\end{lem}

\noindent Proof. Let $\mathcal M_k$ be the $(k+1)\times(k+1)$ second-moment matrix, divided by $T$, of the regressors $(T^{-1/2}\hat u_{t-1},\Delta\hat u_{t-1},\dots,\Delta\hat u_{t-k})$ of the augmented Dickey-Fuller regression, $m_k$ the vector of their moments with $\Delta\hat u_t$, and $\tilde{\mathcal M}_k$ and $\tilde m_k$ the same objects computed from $\tilde u_t$.

Step 1 (the estimated cycle). The difference $\chi_t=\hat u_t-\tilde u_t=-V_t'(\hat\vartheta-\vartheta)-(\hat\mu-\tilde\mu)$ is small. By Lemma \ref{lem:orders}, $\hat\vartheta-\vartheta=O_p(T^{-1/2})$, and the GLS estimates of the constant differ by $(d^{\bar\alpha\prime}d^{\bar\alpha})^{-1}d^{\bar\alpha\prime}V^{\bar\alpha}(\hat\vartheta-\vartheta)=O_p(T^{-1/2})$. In changes, $\Delta\chi_t=-\Delta V_t'(\hat\vartheta-\vartheta)$ for $t\geq2$, so $T^{-1}\sum_t(\Delta\chi_t)^2=O_p(T^{-1})$. In levels, $\max_t|\chi_t|=O_p(1)$ because $\max_t\|x_t\|=O_p(T^{1/2})$, so $T^{-2}\sum_t\chi_t^2=O_p(T^{-1})$.

Step 2 (moments of the augmented Dickey-Fuller regression). By the Cauchy-Schwarz inequality and Step 1, every element of $\mathcal M_k-\tilde{\mathcal M}_k$ and of $m_k-\tilde m_k$ that involves lagged changes is bounded by
\[
2\Big(T^{-1}\sum_t\Delta\tilde u_t^2\Big)^{1/2}\Big(T^{-1}\sum_t\Delta\chi_t^2\Big)^{1/2}+T^{-1}\sum_t\Delta\chi_t^2=O_p(T^{-1/2}),
\]
and the elements that involve the lagged level satisfy the same bound, with $T^{-2}\sum_t\tilde u_t^2=O_p(1)$ and $T^{-2}\sum_t\chi_t^2=O_p(T^{-1})$. The bounds do not depend on the lag indices, so $\max_{k\leq k_T}\|\mathcal M_k-\tilde{\mathcal M}_k\|=O_p(k_TT^{-1/2})$ and $\max_{k\leq k_T}\|m_k-\tilde m_k\|=O_p(k_T^{1/2}T^{-1/2})$.

Step 3 (coefficients and residual variance). The smallest eigenvalue of $\tilde{\mathcal M}_k$ is bounded away from zero uniformly in $k\leq k_T$ with probability tending to one \citep{Berk1974}, and the coefficient vectors $\tilde\beta_{\cdot,k}$ are bounded uniformly in $k$, by the summability of the coefficients of $v_t$. Hence $\max_{k\leq k_T}\|\hat\beta_{\cdot,k}-\tilde\beta_{\cdot,k}\|=O_p(k_TT^{-1/2})$, and the sums of the lag coefficients satisfy
\[
\max_{k\leq k_T}\Big|\sum_{j=1}^k\hat\beta_{j,k}-\sum_{j=1}^k\tilde\beta_{j,k}\Big|\leq k_T^{1/2}\,O_p(k_TT^{-1/2})=O_p\big((k_T^3/T)^{1/2}\big)=o_p(1).
\]
The residual variance is the second moment of $\Delta\hat u_t$ minus $m_k'\mathcal M_k^{-1}m_k$, so by the same bounds it differs from its counterpart for $\tilde u_t$ by $O_p(k_TT^{-1/2})=o_p(1)$, uniformly in $k\leq k_T$. Since $1-\sum_j\tilde\beta_{j,k}$ is bounded away from zero for $\underline k_T\leq k\leq k_T$, the continuous mapping theorem gives the first statement.

Step 4 (the selected order). For $\underline k_T\leq k\leq k_T$, the autoregressive estimator computed from the GLS-demeaned cycle is consistent for $\omega^2$ uniformly in $k$, under the null and under local alternatives, by the results for autoregressions of growing order with a unit or near-unit root in the level term \citep{Berk1974,SaidDickey1984,PerronNg1998}. On the event $\underline k_T\leq\hat k\leq k_T$, $|\hat\omega^2(\hat k)-\omega^2|\leq\max_{k\leq k_T}|\hat\omega^2(k)-\tilde\omega^2(k)|+\max_{\underline k_T\leq k\leq k_T}|\tilde\omega^2(k)-\omega^2|$, and both terms are $o_p(1)$.\footnote{\citet[Section 6.2]{NgPerron2001} note that, under standard asymptotics, the order selected by the modified information criterion grows at a logarithmic rate \citep{HannanDeistler1988}, and their results for the tests are stated for a consistent estimator of the long-run variance. The condition on $\hat k$ is the one under which that consistency holds; with a finite autoregression for $v_t$ it requires only that the criterion does not select fewer lags than the true order.} $\square$

\paragraph{Proof of Theorem \ref{thm:po}.} Write $E_Q=Q_{\bar\alpha}-Q_1$ and $E_A=A_{\bar\alpha}-A_1$ for the changes that quasi-differencing at $\bar\alpha$ instead of $1$ makes to the moments of \eqref{eq:Sdecomp}. By \eqref{eq:Sdecomp},
\[
P_T=\frac{S_u(\bar\alpha)-\bar\alpha S_u(1)}{\hat\omega^2}-\frac{A_{\bar\alpha}'Q_{\bar\alpha}^{-1}A_{\bar\alpha}-\bar\alpha A_1'Q_1^{-1}A_1}{\hat\omega^2}.
\]
The first term is $P^u_T$ up to the estimator of the long-run variance, which is consistent for $\omega^2$ by Lemma \ref{lem:lrv}. For the second, $E_Q,E_A=O_p(1)$, $A_1=O_p(T^{1/2})$ and $T^{-1}Q_1\to_p\Gamma$ by Lemma \ref{lem:orders}. Then $Q_{\bar\alpha}^{-1}-Q_1^{-1}=-Q_1^{-1}E_QQ_{\bar\alpha}^{-1}=O_p(T^{-2})$ and
\[
A_{\bar\alpha}'Q_{\bar\alpha}^{-1}A_{\bar\alpha}-\bar\alpha A_1'Q_1^{-1}A_1=(1-\bar\alpha)A_1'Q_1^{-1}A_1+2E_A'Q_1^{-1}A_1+E_A'Q_1^{-1}E_A+A_{\bar\alpha}'(Q_{\bar\alpha}^{-1}-Q_1^{-1})A_{\bar\alpha},
\]
whose four terms are $O_p(T^{-1})$, $O_p(T^{-1/2})$, $O_p(T^{-1})$ and $O_p(T^{-1})$. Hence $P_T=P^u_T+o_p(1)$ under the null and local alternatives. For (a), a constant satisfies Condition B of \citet{ERS1996}, so $P^u_T\Rightarrow\bar c^2\int_0^1W_c(s)^2ds-\bar c\,W_c(1)^2$ (their Theorem 2), and this limit involves neither $K$, nor $r$, nor the policy terms. For (b), $P^u_T$ is the statistic of the point-optimal test computed from the counterfactual cycle with known weights and policy coefficients. By Theorems 1 and 2 of \citet{ERS1996}, that test is asymptotically equivalent to the Neyman-Pearson test against $c=\bar c$, and its power at $\bar c$ equals the envelope $\Pi(\bar c)$, an upper bound for the power of any test of a unit root in the cycle under Gaussian errors. Since $P_T$ does not use the weights and has the same limit, it belongs to the class of tests that do not know them and attains at $\bar c$ the bound of that class, which cannot exceed $\Pi(\bar c)$. $\square$

\paragraph{Proof of Theorem \ref{thm:diff}.} The decomposition in the proof of Theorem \ref{thm:cycle} uses only $\lambda=\Lambda'b$, so it holds under Assumption \ref{assu:A2d}, and $G_{t+h}-G_{t-1}=(G^0_{t+h}-G^0_{t-1})+\sum_{j\geq0}\tau_j(W_{t+h-j}-W_{t-1-j})$. Multiplying by $u_t$ and applying both parts of Lemma \ref{lem:dl} gives the first statement of (a). For the second, $G_{t+h}-G_{t-1}=\sum_{j=0}^h\Delta G_{t+j}$, and the second part of the proof of Theorem \ref{thm:cycle}, applied to the stationary variable $\Delta G_t$ of $Z^\Delta_t$, gives $\E[\Delta G_{t+j}u_t]/\E[u_t^2]=\gamma^\Delta_j$. For (b), Assumption \ref{assu:A4d}(ii) sets the second term to zero, so $\Gamma_h=\sum_{j=0}^h\tau_j\psi_{h-j}$ with $\psi_0=1$, the system solved in the proof of Theorem \ref{thm:direct}. For (c), Assumption \ref{assu:A2d} with $\lambda\neq\Lambda'b$ gives $Y^1_t(0)-b'Y^0_t=(\lambda-\Lambda'b)'F_t+e_t-b'v_t$, so the cumulated change of the cycle contains the additional term $(\lambda-\Lambda'b)'(F_{t+h}-F_{t-1})$, and the argument of part (a) applies to the rest. $\square$

\paragraph{Proof of Proposition \ref{prop:wchanges}.} Let $\xi_t$ and $\zeta_t$ denote $\Delta Y^0_t$ and $\Delta Y^1_t$ after the projection on the policy terms, and let $\Xi=\Lambda'D^{-1}\Lambda$, the precision of the controls' changes as signals of the common trends. By Assumption \ref{assu:A2d}, $\xi_t=\Lambda\Delta F_t+\Delta v_t$ and $\zeta_t=\lambda'\Delta F_t+\Delta e_t$, so that, under the stated covariances, $\operatorname{Var}(\xi_t)=\Lambda\Sigma_F\Lambda'+D$ and $\operatorname{Cov}(\xi_t,\zeta_t)=\Lambda\Sigma_F\lambda$. The changes are stationary and ergodic by Assumption \ref{assu:A4d}(i), so the sample moments converge and the least-squares weights converge to $w^\ast=\operatorname{Var}(\xi_t)^{-1}\operatorname{Cov}(\xi_t,\zeta_t)$. The identity $(\Lambda\Sigma_F\Lambda'+D)^{-1}\Lambda=D^{-1}\Lambda(I_r+\Sigma_F\Xi)^{-1}$ gives $\Lambda'w^\ast=\Xi(I_r+\Sigma_F\Xi)^{-1}\Sigma_F\lambda=\Xi\Sigma_F(I_r+\Xi\Sigma_F)^{-1}\lambda=\big(I_r-(I_r+\Xi\Sigma_F)^{-1}\big)\lambda$, hence $\lambda-\Lambda'w^\ast=(I_r+\Xi\Sigma_F)^{-1}\lambda$. The cycle built with $w^\ast$ loads on $F_t$ with $\lambda-\Lambda'w^\ast$, and part (c) of Theorem \ref{thm:diff} applies with $b=w^\ast$. With $r=1$, $\Xi\Sigma_F=s$, so $(I_r+\Xi\Sigma_F)^{-1}=1/(1+s)$, and the response of the weighted controls to the innovation is $w^{\ast\prime}\Lambda$ times the response of the trend, which equals $s/(1+s)$ times $\lambda$ times that response, so the added term is $1/s$ times the response of the weighted controls. $\square$

\section{Data sources\label{sec:data}}

\begin{table}[H]
\centering\small
\caption{Data sources.\label{tab:data}}
\begin{tabular}{p{5cm}p{9cm}}
\toprule
Series & Source\\
\midrule
State unemployment rates & Bureau of Labor Statistics, Local Area Unemployment Statistics, seasonally adjusted (FRED series \texttt{TXUR} for Texas and the analogous codes for other states)\\
Alaska and donor-state labour market outcomes & Bureau of Labor Statistics, Local Area Unemployment Statistics (employment-population ratio, unemployment and participation rates) and Current Employment Statistics (payroll employment, leisure and hospitality, retail trade), seasonally adjusted and unadjusted\\
Permanent Fund Dividend & Alaska Department of Revenue, Permanent Fund Dividend Division, Summary of Applications and Payments\\
Oil price and consumer prices & FRED series \texttt{WTISPLC} (West Texas Intermediate) and \texttt{CPIAUCSL} (consumer price index)\\
Weeks of EB by state & Department of Labor trigger notices; replication files of \citet{BooneEtAl2021}\\
Weeks of EUC by state & Replication files of \citet{BooneEtAl2021}, openICPSR project 118525\\
Retail regular gasoline prices & Energy Information Administration, weekly retail gasoline prices for California, the East Coast, Midwest, Gulf Coast and Rocky Mountain regions, Texas and Florida\\
LCFS credit prices & California Air Resources Board, LCFS credit transfer reports, monthly average price\\
LCFS cost per gallon reported by refiners & California Energy Commission, estimated gasoline price breakdown and margins (reports under Senate Bill 1322), May 2024 onward\\
LCFS benchmarks and carbon intensities & California Code of Regulations, title 17, section 95484\\
\bottomrule
\end{tabular}
\end{table}

\bibliographystyle{apalike}
\bibliography{bibliography}

\end{document}
